% Options for packages loaded elsewhere
\PassOptionsToPackage{unicode}{hyperref}
\PassOptionsToPackage{hyphens}{url}
\documentclass[
  ngerman,
]{article}
\usepackage{xcolor}
\usepackage{amsmath,amssymb}
\setcounter{secnumdepth}{-\maxdimen} % remove section numbering
\usepackage{iftex}
\ifPDFTeX
  \usepackage[T1]{fontenc}
  \usepackage[utf8]{inputenc}
  \usepackage{textcomp} % provide euro and other symbols
\else % if luatex or xetex
  \usepackage{unicode-math} % this also loads fontspec
  \defaultfontfeatures{Scale=MatchLowercase}
  \defaultfontfeatures[\rmfamily]{Ligatures=TeX,Scale=1}
\fi
\usepackage{lmodern}
\ifPDFTeX\else
  % xetex/luatex font selection
\fi
% Use upquote if available, for straight quotes in verbatim environments
\IfFileExists{upquote.sty}{\usepackage{upquote}}{}
\IfFileExists{microtype.sty}{% use microtype if available
  \usepackage[]{microtype}
  \UseMicrotypeSet[protrusion]{basicmath} % disable protrusion for tt fonts
}{}
\makeatletter
\@ifundefined{KOMAClassName}{% if non-KOMA class
  \IfFileExists{parskip.sty}{%
    \usepackage{parskip}
  }{% else
    \setlength{\parindent}{0pt}
    \setlength{\parskip}{6pt plus 2pt minus 1pt}}
}{% if KOMA class
  \KOMAoptions{parskip=half}}
\makeatother
\ifLuaTeX
\usepackage[bidi=basic,shorthands=off,]{babel}
\else
\usepackage[bidi=default,shorthands=off,]{babel}
\fi
\ifLuaTeX
  \usepackage{selnolig} % disable illegal ligatures
\fi
\setlength{\emergencystretch}{3em} % prevent overfull lines
\providecommand{\tightlist}{%
  \setlength{\itemsep}{0pt}\setlength{\parskip}{0pt}}
\usepackage{bookmark}
\IfFileExists{xurl.sty}{\usepackage{xurl}}{} % add URL line breaks if available
\urlstyle{same}
\hypersetup{
  pdftitle={Die Zeitlosigkeit der Steine},
  pdfauthor={Paul Koop},
  pdflang={de},
  hidelinks,
  pdfcreator={LaTeX via pandoc}}

\title{Die Zeitlosigkeit der Steine}
\usepackage{etoolbox}
\makeatletter
\providecommand{\subtitle}[1]{% add subtitle to \maketitle
  \apptocmd{\@title}{\par {\large #1 \par}}{}{}
}
\makeatother
\subtitle{Eine Erzählung aus dem Pompeji-Projekt}
\author{Paul Koop}
\date{}

\begin{document}
\maketitle

{
\setcounter{tocdepth}{3}
\tableofcontents
}
\section{Prolog: Die verlangsamte
Zeit}\label{prolog-die-verlangsamte-zeit}

Der Seminarraum lag im dritten Stock, wo die Raumluft nach Ozon,
aufgewärmtem Kunststoff und altem Papier roch. Draußen schnitt der
Septemberwind durch die Straßenschluchten von Boston. Drinnen fiel das
Licht durch die Lamellen der Jalousien, schmale graue Streifen, die in
präzisen Abständen über das Pult und den Dielenboden gezogen waren. Es
wirkte wie ein Barcode auf dem Holz.

Luca saß in der dritten Reihe, die Unterarme flach auf das Kunstharz der
Tischplatte gelegt. Das Tuch seines grauen Pullovers war rau an den
Ellbogen.

An der Stirnwand flimmerte die Projektion. Auf dem Bildschirm lief die
Szene am Bahnsteig -- Jaco Van Dormaels \emph{Mr. Nobody}, Kadenz für
Kadenz heruntergedimmt auf eine beinahe unerträgliche Langsamkeit. Der
neunjährige Nemo stand im geometrischen Schnittpunkt zweier Blicke. Die
Mutter mit dem Koffer auf dem Bahnsteig; der Vater neben dem Jungen, das
Gesicht gezeichnet von einer stummen, erschütterten Verzweiflung, die
Ränder der Augen gerötet. Die Kamera hielt beide Eltern zugleich im
Kader, während die Unschärfe wechselte. Nemo blickte von der Mutter zum
Vater und zurück, gefangen in einer räumlichen Anordnung, aus der er
keinen der beiden ansehen konnte, ohne den anderen im Vordergrund zu
spüren.

„Willst du mitkommen oder beim Vater bleiben?{\kern0pt}``, fragte die
Stimme der Mutter aus den Lautsprechern. Die Worte klangen flach, fast
künstlich im gedämpften Raum.

Der Zug fuhr an. Die Zeit für den Entschluss lief ab. Das Bild
verzweigte sich im Schnitt: Nemos innere Lähmung schlug um in äußere
Hektik. In einer Einstellungsfolge rannte der Junge dem fahrenden Waggon
nach, erreichte die Hand der Mutter. In der nächsten Kameraeinstellung
löste sich beim Anlaufen sein Schnürsenkel. Ein völlig kontingenter
Moment -- das Stolpern, der fallende Schuh auf den Betonplatten, das
Abreißen des Kontakts. Der Zufall entschied über die Existenzachse,
bevor die Wahl überhaupt vollzogen war. \emph{„Daddy, is it my
fault?{\kern0pt}``}, fragte die Kinderstimme im Rauschen des
Videoband-Audios, worauf der Vater abwinkte: \emph{„Nein, es ist nicht
deine Schuld.``}

„Wir sprechen hier nicht von einer einfachen Dramaturgie der Wahl``,
sagte der Dozent. Er stand seitlich des Strahlers, sodass das Licht sein
Gesicht halbierte. Die Brille reflektierte das Flimmern des Bahnsteigs.
„Van Dormael zeigt den Knotenpunkt. Die Entscheidung ist kein souveräner
Willensakt, sondern ein Zusammenspiel aus Impuls, verstreichender Zeit
und purer Kontingenz -- einem offenen Schnürsenkel. Wenn das InSIM-Labor
in Mailand mit solchen Matrizen arbeitet, tun sie genau das. Über eine
Quanten-Schnittstelle füttern sie das System mit probabilistischen
Pfaden. Die Viele-Welten-Interpretation wird dort zur Rechenarchitektur:
Der Junge entscheidet sich nicht bloß für A oder B. Er verzweigt sich in
alle Zustände gleichzeitig.``

Er trat einen Schritt zur Seite. Der Staub im Lichtkegel tanzte in
Mikrobewegungen, winzige Partikel, die stumm kollidierten und abperlten.

„Teilhard de Chardin sah im Omega-Punkt die ultimative Verdichtung von
Bewusstsein``, fuhr der Mann fort, während er die Kreide zwischen zwei
Fingern drehte, bis weißer Staub seine Knöchel bedeckte. „Ein
Zusammenstürzen der Zeitachse auf ein einziges, alles umfassendes
Zentrum. Aber was passiert, wenn die Rechenleistung eines Systems --
nennen wir es das ARS-Modell -- diese Singularität auflöst und jede
Möglichkeit, jedes Stolpern am Bahnsteig, simultan real hält? Wo bleibt
die Schuld, wo bleibt der Wille, wenn jede Kontingenz eine neue Welt
erzeugt?{\kern0pt}``

Luca sah die Unschärfe auf dem Monitor. Seine Handschuhe lagen neben dem
Notizbuch, das aufgeschlagen, aber unbeschrieben war. Die Seite zeigte
das Raster eines linearen Liniensystems.

„Für die nächste Sitzung``, sagte der Dozent und trat an den Tisch, um
den Laptop einzuklappen. Das Scharnier knarrte leise. „Lesen Sie die
Arbeiten von Michael Phillips über die empirische Bestimmung von
Dialoggrammatiken. Und vergleichen Sie seine Formalisierung des
Nicht-Sprechens mit der Rahmenhandlung von DeLillos \emph{Point Omega}.
Achten Sie auf die Stille im Wüstenraum. Auf die Unmöglichkeit, ein
Verschwinden zu protokollieren, wenn der Beobachter selbst in den
Verzweigungen gefangen ist.``

In den Reihen um Luca herum begann das mechanische Rauschen des
Aufbruchs. Reißverschlüsse schlossen sich, schwere Stoffe rieben
aneinander, Stuhlbeine schrammten über das Linoleum. Die anderen
bewegten sich in der geläufigen Taktung des Nachmittags.

Luca blieb sitzen.

Der Dozent klemmte sich die Ledermappe unter den Arm. Er sah kurz zu
Luca hinab, ohne die Hand vom Schreibtischrand zu nehmen.

„Sie suchen nach einer Logik in den Abweichungen, Luca``, sagte er. Die
Stimme war leise, frei von akademischer Geste. „Aber Sprachmodelle und
stochastische Prozesse liefern keine Urheber. Sie liefern nur
Kohärenzen.``

„Er antwortet nicht auf die Mails der Fakultät``, sagte Luca.

„Phillips hat das Jesuiten-Institut im Frühjahr verlassen``, entgegnete
der Dozent, während er zur Tür blickte. „Er ist jenseits der
Datenkanäle. Physiker von seinem Schlag neigen im Alter zu einer
eigentümlichen Form von Konservierung. Sie wollen die Dinge erstarren
sehen. Wie die Abgüsse unter der Asche in der Campania.``

Luca hob den Kopf nicht. „Er wartet.``

Der Dozent sah ihn noch einen Moment lang an, die Lippen leicht
zusammengedrückt, dann wandte er sich ab. Seine Ledersohlen erzeugten
ein trockenes Quietschen auf dem Gang, bis das Geräusch in der Geometrie
des Treppenhauses verhallte.

Luca stand auf. Der Raum war kühl geworden. Er schob den Stuhl nicht
zurück.

Auf dem Flur lag die Nachmittagssonne tief über der Kulisse von Boston.
Unter den Glasfronten der Bankentürme schoben sich die Fahrzeugkolonnen
im Stop-and-Go durch die Schluchten, ein träger Fluss aus Blech und
Bremslichtern, dessen Lärm durch die Doppelverglasung zu einem dunklen,
mechanischen Summen gedämpft wurde.

Er griff in die Brusttasche seiner Jacke. Das Papier des ungeöffneten
Bogen-Spiegels war rau. Keine Frankierung. Kein Stempel eines Postamtes.
Nur die maschinell gesetzten Typen auf dem vergilbten Dokumentenpapier,
deren Tinte an den Rändern leicht zerflossen war, wie die Ausgabe eines
Thermodruckers:

\emph{Er steht im Versuchsfeld. An der Grenze zur Senke. Wenn du ihn
suchst, sprich nicht von der Zeit. Er rechnet nicht mehr.}

Luca ließ das Papier wieder in die Tasche gleiten. Die Luft, die durch
den Türspalt des Notausgangs strömte, schmeckte nach trockenem Ruß und
Staub. Er ging die Treppe hinunter, Stufe für Stufe, im gleichmäßigen
Rhythmus eines Prozesses, der längst gestartet war und sich nicht mehr
abbrechen ließ.

\section{\texorpdfstring{\textbf{Kapitel 1: Der
Brief}}{Kapitel 1: Der Brief}}\label{kapitel-1-der-brief}

Die Wohnung im vierten Stock lag tief im Schatten der gegenüberliegenden
Backsteinfronten. Unter dem Fenster verlief die schmale Ashburton Place.
Das Geräusch der Busse war ein rhythmisches Entweichen von Pressluft,
das von Zeit zu Zeit die Scheiben vibrieren ließ. Ein kühles,
metallisches Zittern im Rahmen.

Luca saß am Küchentisch, der aus geschichteter Spanplatte bestand. Das
Papier lag zwischen einer Schale mit verblasstem Trockenobst und der
Porzellantasse. Es war zweimal gefaltet, flachgepresst von der Türkante,
unter der man es am Morgen durchgeschoben hatte. Keine Frankierung. Ein
ungleichmäßiges Blatt aus einem Nadeldrucker, die Ränder feingezackt
durch abgerissenes Endlospapier.

\emph{Er steht im Versuchsfeld. An der Grenze zur Senke. Wenn du ihn
suchst, sprich nicht von der Zeit. Er rechnet nicht mehr.}

Die Schrift zeigte keine Neigung. Es war der Ausdruck einer alten
Terminalschnittstelle, kühle Matrizen ohne personellen Duktus. Luca fuhr
mit der Kuppe des Zeigefingers über die Zeile. Die winzigen Punkte des
Farbbandes lagen tief in den Fasern des Papiers.

Draußen streifte ein Streifenwagen im Schneckentempo die Bordsteinkante.
Ein Hund schnüffelte an einem Hydranten, stoppte, ging weiter. Das
Geräusch des Motors verflachte in Richtung Cambridge Street.

In der Küche schabte Porzellan auf der Granitimitat-Arbeitsplatte. Seine
Mutter füllte die Kanne. Der Wasserkocher gab ein scharfes, metallisches
Zischen von sich, das schlagartig abbrach, sobald sie den Schalter
umlegte. Sie trat in den Türrahmen, ein knappes Tableau aus
ausgewaschenem Baumwollstoff und grauen Haarnadeln. Sie trug die Tassen
nicht auf einem Tablett, sondern hielt sie fest an den Henkeln
umschlungen, die Knöchel weiß hervortretend.

Sie stellte die Tassen ab. Der Wasserdampf stieg geradlinig auf, bevor
er im Durchzug der Raumluft zerfiel.

„Die Schnittstelle in Mailand hat zwei Abfragen registriert``, sagte
Luca, ohne aufzusehen. „Um vier Uhr morgens. Datenströme über den Server
der Gregoriana.``

Sie setzte sich auf die Stuhlkante. Ihre Ellbogen stützten sich nicht
ab; sie hielt den Oberkörper in einer unnatürlichen, gelernten Distanz
zur Tischkante.

„Das sind automatisierte Abgleiche``, sagte sie. Die Stimme war dünn,
frei von Frequenz. „InSIM testet die Dialoggrammatiken. Das macht das
System von selbst, wenn die Knotenpunkte inaktiv bleiben.``

„Das ist keine Dialoggrammatik``, sagte Luca und tippte auf das Papier.
„Das ist eine Adresse.``

Sie sah nicht auf den Bogen. Ihre Augen fixierten das blasse
Lichtmuster, das sich brechend im feuchten Boden ihrer Tasse spiegelte.
„Phillips hat keine Adresse mehr. Wer auch immer den Prompter bedient
hat, nutzt seine alten Parameter. Man lässt Algorithmen nicht allein
laufen, wenn sie auf Quantenschnittstellen zugreifen. Sie fangen an,
Ausgänge zu erfinden.``

„Er ist in der Sonora-Senke.``

„Dort steht nichts als Staub und die Relikte einer Funkmessstation``,
antwortete sie. „Ein Konservierungsexperiment aus den Siebzigern.
Trockenheit, die den Zerfall aufschiebt. Sie haben dort früher
Gesteinsproben aus der Campania gelagert, um den Abrieb unter
simulierter Sonneneinstrahlung zu messen. Die zeitlose Substanz.``

Luca sah sie an. Die Haut unter ihren Augen war transparent, fein
durchzogen von bläulichen Äderchen. Sie sprach von der Wüste wie von
einer Laboranordnung, einem isolierten Versuchsraum für Zerfallsraten.

„Er ist dein Vater nicht im biologischen Sinne einer Entscheidung``,
sagte sie leise. „Er war die Versuchsbedingung. Die Jesuiten nannten es
eine Beobachter-Funktion.``

Luca stand auf. Der Linoleumboden war kalt unter seinen Socken. Er ging
in das schmale Zimmer neben der Küche, ziehen, greifen, falten. Der
Stoff der Reisejacke roch nach dem muffigen Schrankfutter. Er schob zwei
Hemden, ein Notizbuch ohne Rasterung und die Reserveakkus für das
Diktiergerät in den Synthetikbeutel. Seine Bewegungen folgten einem
trockenen, funktionalen Ablauf, wie das Abarbeiten einer Checkliste vor
einem Systemstart.

Als er zurück in die Küche trat, lag ein verblasstes Farbfoto neben der
ungeöffneten Tasse. Es war ein Kleinbildabzug aus den frühen Neunzigern,
die Ränder durch Lichteinfall gelblich verfärbt.

Das Bild zeigte einen Mann in einem dunklen Kragenhemd vor einer
staubigen Ziegelwand im Campo Santo Teutonico. Das Gesicht war schmal,
die Schläfen eingefallen, die Augen unruhig in den Schatten der Höhlen
zurückgezogen -- der Blick eines Mannes, der gelernt hatte, Messreihen
zu lesen, statt Gesichter zu dekodieren.

„Er glaubte an die Verschränkung``, sagte seine Mutter vom Fenster aus.
Sie sah hinab auf die Fahrbahn der Ashburton Place. „Er sagte immer,
wenn zwei Zustände lang genug gekoppelt bleiben, gibt es kein
Getrenntsein mehr. Selbst dann nicht, wenn einer von beiden erlischt.
Man sucht dann nicht nach einer Person. Man sucht nach dem Echo des
Zerfalls.``

Luca steckte das Foto zwischen die Seiten des Notizbuchs. Der Bogen mit
dem Ausdruck wanderte in die Innentasche seiner Jacke, direkt an den
Rippenbogen.

„Ich nehme um 14:10 Uhr die Red Line``, sagte er. „An South Station
steige ich um in die Silver Line zum Logan Airport.``

Sie drehte sich nicht um. Ihre Hand lag flach auf der kalten Scheibe,
hinter der das Blaulicht eines entfernten Einsatzwagens stumm über die
Fassaden der gegenüberliegenden Gebäude glitt. „Wenn du dort ankommst,
wirst du keine Antworten finden, Luca. Nur die Reste der Daten, die ARS
noch nicht bereinigt hat.``

Er zog den Reißverschluss der Tasche zu. Das Geräusch der
Kunststoffzähne war das einzige Signal im Raum. Er öffnete die Tür zum
Treppenhaus. Der Geruch von feuchtem Beton und altem Holz trat ihm
entgegen.

„Schalt das Terminal aus``, sagte er.

„Es läuft über das Netz von InSIM``, antwortete sie leise in den leeren
Raum hinein. „Man kann ein Quantensystem nicht ausschalten, indem man
den Stecker zieht. Es bleibt im Überlagerungszustand.``

Die Tür fiel ins Schloss. Der Riegel schnappte mit einem trockenen,
metallischen Klick ein.

Auf der Treppe verlangsamte Luca den Schritt nicht. Mit jeder Stufe, die
er hinabstieg, wurde das Dröhnen der Stadt draußen deutlicher -- das
abgedämpfte Rollen der Reifenspuren auf nassem Asphalt, das Pfeifen des
Windes im Belüftungsschacht der U-Bahn-Station an der Park Street.

Im Waggon der Red Line saß er der eigenen Spiegelung in der geschwärzten
Fensterscheibe gegenüber, während der Zug durch den Tunnel unter dem
Charles River schnitt. Das Gesicht im Glas war unscharf, überlagert von
den Signalleuchten im Tunnel, die rhythmisch in die Dunkelheit
schnitten. Gelb. Rot. Aus. Ein Stroboskop aus Möglichkeiten, das ihn
unverändert weiter nach Osten schob.

\section{\texorpdfstring{\textbf{Kapitel 2: Die
Reise}}{Kapitel 2: Die Reise}}\label{kapitel-2-die-reise}

Sechs Stunden im Rumpf einer Boeing 777, isoliert über der
Kontinentalplatte. Die Kabinenluft war auf zweiundzwanzig Grad
heruntergekühlt, trocken und mit dem Geruch von künstlicher Zitrone,
Kaffee und einem schwachen Hauch von Kerosin gesättigt.

Luca saß am Fenster auf Platz 14A. Vor ihm auf dem Klapptisch stand die
Plastiktasse mit dem Filtermilchkaffee. Eine hauchdünne, dunkle Haut
hatte sich auf der Oberfläche gebildet, die bei jeder leichten
Kurvenlage des Flugzeugs erschütterungsfrei an den Becherrand stieß. Er
rührte sie nicht an.

Draußen lag die Wolkendecke flach und geschlossen unter der Tragfläche,
ein unendliches Feld aus strukturlosem Weiß. Es sah aus wie die gipserne
Oberfläche der Geländemodelle im InSIM-Labor -- eine glatte, noch nicht
durch Berechnungen aufgelöste Matrix.

Er zog das Farbfoto aus dem Notizbuch. Der Mann vor der Ziegelwand im
Campo Santo Teutonico. Im scharfen Licht der Leseleuchte traten die
Kanten des Papierabzugs noch deutlicher hervor. Phillips' Augen waren
nicht auf die Linse gerichtet, sondern leicht seitlich versetzt, als
wartete er auf das Signal einer Messapparatur außerhalb des
Bildausschnitts. Die Gesetze der Quantenmechanik: Ein Zustand, der erst
durch den Akt der Beobachtung determiniert wird. Was geschah mit der
Wellenfunktion, wenn der Beobachter die Zelle verließ?

Die Flugbegleiterin schob den Servierwagen durch den Gang. Das Radlager
der Rollen erzeugte ein hohes, mechanisches Summen auf dem Teppichboden.

„Noch etwas Wasser für Sie?{\kern0pt}``, fragte sie. Das Lächeln war
perfekt kalibriert, eine geschulte Dialoggrammatik für Zehntausend Meter
Höhe.

„Nein``, sagte er.

„Wir beginnen in Kürze mit dem Sinkflug auf Los Angeles.``

Er sah wieder durch die doppelte Acrylscheibe. Die dichte Schicht löste
sich auf. Die Wolken zerrissen in flache, graue Fetzen, die über die
Topographie Kaliforniens schwebten. Darunter erschien das Land: karg,
erdfarben, durchzogen von trockenen Flussbetten, die wie verblasste
Schaltungsmuster in die Kruste gezeichnet waren.

Dann die Agglomeration. Los Angeles war kein zusammenhängender Ort,
sondern ein Raster aus Freeways, Parkdecks und flachen Gewerbebauten,
das sich bis zum Dunstschleier des Pazifiks erstreckte. Eine schier
unendliche Redundanz aus Beton und Asphalt, in der Millionen Bewegungen
simultan abliefen, ohne einander zu berühren.

LAX roch nach Ozon, Flugbenzin und dem feuchten Gummi der Gepäckbänder.
Die Transitpassagiere schoben sich stumm unter den Anzeigetafeln
hindurch, deren Pixelanzeigen im Sekundentakt neue Ankunftszeiten
ausgaben.

Am Schalter der Autovermietung unterschrieb Luca den Standardvertrag.
Der Stift war an einer Spiralschnur aus schwarzem Kunststoff befestigt.

„Kompaktklasse, grauer Sedan, Deck 3, Stellplatz 412``, sagte die
Angestellte, ohne den Blick vom Monitor zu heben. Ihre Finger glitten
routiniert über die Tastatur.

Im Parkhaus stand der Wagen unter den Neonröhren. Das Metall des Dachs
war kalt. Der Motor sprang sofort an -- ein leises, vibrationsarmes
Brummen des Vierzylinders.

Er steuerte den Wagen auf die Interstate 10 East. Der Verkehr auf den
sechs Spuren floss mit sechzig, siebzig Meilen pro Stunde, eine dichte
Kolonne aus Chrom, Lack und verdunkelten Scheiben. Neben ihm zog die
Peripherie vorbei: Plakatwände für Logistikunternehmen,
Fast-Food-Schilder auf Stahlmasten, automatisierte Waschanlagen. Alles
wirkte wie vorgefertigt, aufgestellt in einem Versuchsfeld für
funktionale Architektur.

Nach zwei Stunden veränderten sich die Radien der Straße. Die Wohnblöcke
wichen flachen Lagerschuppen, dann folgten die Abzweigungen zu den
Fertigsiedlungen der Wüstenvororte. Schließlich hörten auch die Zäune
auf.

Die Straße wurde zu einem grauen Band, das zwischen den letzten
Ausläufern der Berge in die Colorado Desert schnitt. Die Sonne stand
tief im Rückspiegel und warf den schmalen Schatten des Sedans weit vor
das Fahrzeug. Die Landschaft verkümmerte zu verbrannten Büschen, hellen
Kalksteinen und feinem, rotem Staub, der vom Wind quer über den Asphalt
getrieben wurde.

Er schaltete das Radio ein. Rauschen auf allen Frequenzen, unterbrochen
von den Bruchstücken eines Predigers auf Mittelwelle, dessen Stimme im
statischen Knistern unterging. \emph{„...the point where time ceases to
measure...``} Er drehte den Lautstärkeregler nach links, bis der
Schalter klickte.

Die Stille im Fahrzeugraum war absolut, nur überlagert vom gleichmäßigen
Abrollgeräusch der Reifen auf dem rauen Belag.

Er verließ die Interstate und folgte der schmalen Straße nach Südosten.
An einer Abzweigung ohne Ampel stand die einzelne Zapfsäule einer alten
Tankstelle vor einer Baracke aus Wellblech. Der Wind ließ ein
verrostetes Werbeschild an zwei Ketten schlagen.

Luca hielt an. Als er die Fahrertür öffnete, schlug ihm die Hitze
entgegen -- trocken, scharf, durchsetzt mit dem Geruch von verbranntem
Harz und Staub. Es war eine Hitze, die nicht wärmte, sondern
austrocknete, ganz ähnlich der kalten Trockenheit in den klimatisierten
Depots von Pompeji, in denen man die organischen Reste unter
Luftabschluss hielt.

Er ließ den Treibstoff in den Tank laufen. Das Zählwerk an der Zapfsäule
klackerte mechanisch.

Im Verkaufsraum saß ein Mann hinter einem Schutzglas aus Acryl. Seine
Haut war tief gefurcht, das Hemd an den Schultern verwaschen. Auf dem
Tresen lag eine vergilbte Faltkarte des San Bernardino County.

Luca legte zwei Plastikflaschen stilles Wasser und ein Sandwich in Folie
auf das Band.

Der Mann scannte die Barcodes. Der Piepton war schrill im engen Raum.
„Die Nebenstraße ist weiter oben gesperrt``, sagte er. Die Stimme klang
rau, als benutze er sie selten. „Bauarbeiten an der alten Brücke.``

„Ich muss weiter nach Südosten``, sagte Luca. „In die Senke.``

Der Mann sah ihn durch das leicht getönte Acrylglas an. Die Augen waren
schmal, trüb vom grauen Star. „Da ist nichts mehr. Die alten
Messstationen haben sie vor zehn Jahren abgeschaltet. Die Army nutzt das
Terrain nicht einmal mehr für Zielübungen.``

„Die Straße führt durch``, sagte Luca.

„Die Straße endet an der alten Salzpfanne``, entgegnete der Mann und
schob das Wechselgeld unter der Scheibe durch. „Danach gibt es nur noch
Fahrspuren im Sand. Wenn der Motor ausfällt, bleibt die Wärme im Metall.
Bis sie abkühlt. Mehr ist da nicht.``

Luca nahm die Flaschen. „Das reicht.``

Er stieg wieder ein. Die Tür schloss mit einem dumpfen Schlag.

Die Sonne versank hinter den Bergketten im Westen und hinterließ einen
breiten, violetten Streifen am Horizont, der rasch in tiefes Schwarz
überging. Die Scheinwerfer des Sedans schnitten zwei weiße Lichtkegel in
die Dunkelheit, in denen Milliarden Staubpartikel tanzten. Kein
Gegenverkehr mehr. Keine Masten. Nur die gelbe Mittellinie, die im
Rhythmus der Scheinwerfer unter dem Kühler verschwand.

Er hielt am Rand einer geschotterten Ausweichstelle. Schaltete die
Zündung aus. Die Scheinwerfer erloschen.

Sofort brach die Wüstennacht über den Wagen herein. Der Motor knisterte
leise im Abkühlen, Metall, das sich im Schrumpfen zusammenzog.

Luca stieg aus und lehnte sich an die Motorhaube. Der Himmel über ihm
war klar, dicht besetzt mit kalten, feinen Lichtpunkten. Kein
Zivilisationsleuchten mehr am Rand. Die Luft war schlagartig abgekühlt.

Er zog den Brief aus der Tasche. Im schwachen Licht des Displays seines
Mobiltelefons -- kein Empfang, keine Balken in der Statusleiste -- las
er die Zeilen nicht mehr, er spürte nur die Struktur des Papiers.

\emph{Er rechnet nicht mehr.}

Das war die Formulierung von ARS: ein Zustand, in dem ein
Quantenprozessor aufgehört hatte, neue Iterationen zu erzeugen, weil
alle Zustände kollabiert waren.

Gegen zweiundzwanzig Uhr erreichte er ein Motel am Rand einer
aufgelassenen Bergbausiedlung. Drei Leuchtröhren des Neonschilds summten
in der Kälte: \emph{M O T E L}.

Der Schlüssel war aus Messing, befestigt an einer rautenförmigen
Plastikmarke mit der Nummer 7. Das Zimmer roch nach Linoleum,
Reinigungsmittel und trockenem Holz. Auf dem Nachttisch stand ein
Telefon mit Wählscheibe, das Kabel aus der Wand gezogen.

Luca legte sich unentkleidet auf die Decke. Das Licht der einzelnen
Straßenlaterne vor dem Fenster warf das Streifenmuster der Jalousie an
die Decke -- genau wie im Seminarraum in Boston.

Er schloss die Augen nicht. Er lauschte auf das Geräusch des Windes, der
den feinen Wüstenstaub gegen die Scheiben blies. Es klang wie das
Rauschen eines Datenstroms, der langsam im Rauschen des Hintergrunds
erlischt.

\section{\texorpdfstring{\hfill\break
}{ }}\label{section}

\section{\texorpdfstring{\textbf{Kapitel 3: Die
Ankunft}}{Kapitel 3: Die Ankunft}}\label{kapitel-3-die-ankunft}

Der Asphalt fraß sich am Rand der Senke fest, zerbröckelte in
unregelmäßigen Kalksteinstücken und ging schließlich in eine Piste aus
verdichtetem, ockerfarbenem Sediment über. Der Sedan verlangsamte das
Tempo. Der Unterboden setzte einmal hart auf einem auswaschenden
Gesteinsgrat auf -- ein trockenes, metallisches Schlagen, das durch die
Lenksäule bis in Lucas Handgelenke schoss.

Der Staub zog in einer dichten, feinkörnigen Fahne hinter dem Wagen her.
Er setzte sich auf die Heckscheibe, kroch durch die Dichtungen der
Luftschlitze und legte sich wie eine mehlige Schicht auf das
Armaturenbrett. Es schmeckte nach verbranntem Kalk und ausgetrockneter
Tonerde, ein feiner, alkalischer Film auf der Zunge, der an die
Konservierungspulver in den Archäologiedepots der Campania erinnerte.

Am Abhang des Hügels stand die Konstruktion.

Sie war schmal, auf flachen Betonfundamenten errichtet, das Holz von der
ultravioletten Strahlung verbrannt und zu einem matten, silbrigen Grau
verwittert. Die Bretter wiesen tiefe Längsrissmuster auf, aus denen
trockenes Harz getreten war. Kein Dachrinnenrohr, kein Maschendraht. Nur
die rechteckige Geometrie einer abgelegenen Feldstation, deren
Fensterrahmen vom feinen Wüstensand blankgeschmirgelt worden waren.

Auf der schmalen Veranda saß Phillips.

Er lag nicht im Liegestuhl, er lehnte nicht an der Fassade. Er saß auf
einem einfachen Holzstuhl, die Knie rechtwinklig gebeugt, die Füße flach
auf den Dielen. Seine Hände ruhten auf den Oberschenkeln -- lange,
knochige Finger mit verhornten Nägeln, deren Adern wie blaue Kabel unter
der pergamentartigen Haut hervortraten. Er trug ein ausgewaschenes,
dunkles Kragenhemd, dessen oberster Knopf fehlte.

Luca stellte den Motor ab. Die schlagartige Abwesenheit der
Zylinderbewegung erzeugte eine physische Leere im Raum. Der Staub sank
langsam auf die Motorhaube; das heiße Metall tickte im Sekundentakt der
Abkühlung.

Luca stieg aus. Der Sand knirschte trocken unter den Gummisohlen seiner
Sneaker. Die Luft hier unten war anders als an der Zapfsäule --
stehender, schwerer, durchdrungen von der Hitze, die das kalkhaltige
Gestein den ganzen Tag über gespeichert hatte. Sie roch nach Nichts,
nach thermisch gereinigtem Raum.

Er trat an die Stufen der Veranda. Phillips bewegte den Kopf nicht. Sein
Blick war flach über die Geländemulde gerichtet, dorthin, wo die
Überreste der alten Messmasten als rostrote Nadeln im Flimmern der
Horizontlinie standen.

„Die Schnittstelle in Mailand hat Ihre Parameter nicht gelöscht, Dr.
Phillips``, sagte Luca. Die eigene Stimme klang im Freien erstaunlich
flach, als würde die offene Wüste die Stimme sofort verschlucken. „Das
InSIM-System führt Ihre Dialoggrammatiken weiter. Es erzeugt neue
Schleifen.``

Phillips schwieg. Ein Insekt zog in einer trägen Kreisbahn um das
verblichene Holz des Geländers, stumm und ohne Summen.

„Die Fakultät in Boston nimmt an, Sie seien in Rom``, fuhr Luca fort. Er
griff in die Jackentasche und spürte die harten Kanten des Fotos. „Aber
der Bogen aus dem Nadeldrucker kam ohne IP-Header. Er nutzte den alten
Account der Gregoriana.``

Phillips schloss für die Dauer eines Atemzugs die Augen. Seine Lidränder
waren trocken, gezeichnet von feinsten Staubadern. Das Gesicht war
schmaler geworden als auf dem Bild vor der Ziegelwand im Campo Santo --
das Skelett zeichnete sich unter der gegerbten Haut ab wie ein
freigelegtes Artefakt.

„Teilhard schrieb von der Zuspitzung``, sagte Luca, die Schritte bedacht
auf der untersten Stufe ansetzend. „Aber das ARS-Modell rechnet nicht
auf eine Zuspitzung hin. Es verzweigt. Es lässt die Möglichkeiten im
Raum stehen, ohne eine zu tilgen. Was ist das für ein Ort, Phillips? Ein
Beobachterposten?{\kern0pt}``

Das Licht veränderte seine Neigung. Die Schatten der Dachträger schoben
sich millimeterweise über das Nadelholz der Veranda, scharfe, schwarze
Balken, die das Deck in mathematische Segmente unterteilten.

Luca bewegte sich nicht. Das Trikot-Gewebe seines Hemdes klebte am
Rücken. Die Hitze war kein Gefühl, sondern ein Druck auf dem Brustbein,
eine physikalische Einschränkung des Atems. Er fixierte die rechte Hand
des alten Mannes. An der Kante seines Handgelenks verlief eine kleine,
helle Narbe -- die Schnittspur einer alten Glaspipette aus den
Physikalischen Laboren.

Nach einer Stunde -- die Sonne hatte die Kuppe der westlichen
Wüstenkette erreicht und tauchte den Senkenboden in ein hartes, rotes
Licht -- stand Phillips auf. Seine Bewegungen waren mechanisch, ohne die
Hast alter Menschen. Er griff nach der blechernen Kanne auf dem Tisch
hinter sich, füllte ein dickwandiges Wasserglas und schob es an die
Kante des Geländers.

Das Glas war beschlagen, feine Kondensatperlen bildeten sich auf der
Außenseite und liefen durch den Wüstenstaub herab.

Phillips sah ihn erstmals an. Die Augen waren nicht trüb, sondern von
einem blassen, fast durchsichtigen Grau, im Zentrum verengt zu zwei
winzigen, schwarzen Punkten.

„Das Wasser kommt aus der tiefen Bohrung``, sagte Phillips. Die Stimme
war kein Flüstern, sondern ein trockenes, scharfes Geräusch, wie das
Reiben zweier Kalksteine aneinander. „Achtzig Meter. Vor dem Ausbruch
gebunden.``

Luca nahm das Glas. Das Wasser war kalt, schmeckte metallisch nach Eisen
und Natrium, brannte leicht an den ausgetrockneten Lippen.

„Sie haben die Quantenschnittstelle nicht abgeschaltet``, sagte Luca.

Phillips wandte den Kopf wieder ab. Sein Profil schnitt scharf gegen den
glühenden Horizont. „Wenn du nach Sätzen suchst, bist du im falschen
Versuchsfeld. Die Dialoge sind abgeschlossen. Wir dokumentieren hier nur
noch das Erstarren.``

„Das Buch des Systems ist nicht leer``, sagte Luca.

„Das System rechnet die Reste``, antwortete Phillips leise. „Es erzeugt
Zustände, weil niemand da ist, der die Messung beendet. Wenn du bleiben
willst, lass die Fragen im Wagen. Der Staub braucht keine Grammatik.``

Er drehte sich um und trat durch die Fliegengittertür. Das Netz schlug
mit einem leisen, doppelten Klacken gegen den Rahmen.

Luca blieb auf den Stufen sitzen. Das Holz unter seinen Oberschenkeln
kühlte mit dem Verlöschen des Tageslichts schlagartig ab. Draußen in der
Senke verblassten die Schatten der Messmasten zu dünnen, grauen
Strichen, bis die Dunkelheit das Becken wie flüssige Tinte ausfüllte.
Kein Zivilisationsgeräusch. Kein Flugzeugnetz am Nachthimmel. Nur das
feine, mechanische Summen des Windes, der den Sand gegen die Fundamente
trieb.

Später trat Luca in den Innenraum.

Die Hütte roch nach Paraffin, altem Papier und trockenem Zedernholz. Ein
einziges Petroleumlicht brannte auf dem Holztisch und warf einen engen,
gelben Lichtkreis auf die Dielen.

Im Regal an der Rückwand standen keine populären Taschenbücher. Dort
reihten sich englische und italienische Forschungsbände in abgewetzten
Leineneinbänden: Abhandlungen über stochastische Prozesse, Manuskripte
zur Phänomenologie der Zeit, alte Berichte der Päpstlichen Akademie der
Wissenschaften. Dazwischen lagen verblasste Karteikarten, dicht
beschrieben mit mathematischen Formeln und altgriechischen
Textfragmenten -- Phillips' eigene Vorarbeiten zu den Dialoggrammatiken
aus den Neunzigern.

An der Wand über dem feldbettähnlichen Gestell hing kein Kruzifix aus
poliertem Holz, sondern ein einfaches, gusseisernes Kreuz ohne Korpus,
dessen Oberfläche korrodiert war. Die Arme des Eisenkreuzes waren leicht
asymmetrisch, wie ein verlassenes Messergebnis.

Phillips lag auf dem Schlafsack, die Arme über der Brust gekreuzt, die
Augen geschlossen. Sein Atem war so flach, dass sich der Stoff seines
Hemdes kaum hob.

Luca setzte sich auf die freie Pritsche an der gegenüberliegenden Wand.
Er zog das Notizbuch heraus, legte es ungeöffnet auf das hölzerne
Gestell neben seinem Kopf und starrte an die Decke, wo das Licht der
Funzel das Geäst der Balkenstrukturen abbildete.

Die Wüste draußen war kein Ort mehr. Sie war ein Resonanzraum für ein
System, das vergessen hatte, wie man kollabiert. Luca schloss die Augen
und lauschte auf das trockene Ticken der Deckenbalken, die im Abkühlen
die Wärme der Sonne abgaben.

\section{\texorpdfstring{\textbf{Kapitel 4: Das
Schweigen}}{Kapitel 4: Das Schweigen}}\label{kapitel-4-das-schweigen}

Die Entschleunigung vollzog sich nicht als plötzlicher Bruch, sondern
als schleichender Entzug von Impulsen. Die Rhythmen der Tage schrumpften
auf wenige, starr wiederkehrende Sequenzen -- eine Messreihe ohne
Ausreißer.

Das Licht schob sich morgens wie ein graues Lineal durch die Lücken der
Holzverschallung. Phillips saß bereits auf der Veranda, die Unterarme
flach auf den Knien, die Schultern leicht nach vorne gezogen. Er trug
dieselbe ausgewaschene Baumwollhose, deren Saum vom kalkhaltigen Staub
weiß verkrustet war.

Luca trat hinter ihm auf das Deck. Die Dielen gaben keinen Laut von
sich. Er setzte sich auf die oberste Stufe, den Rücken an den Pfosten
gelehnt.

In der Mündung der Senke lag der Dunst flach über den Basaltblöcken. Das
Licht wirkte noch ungefiltert, kalt und brechend hart.

„Die Schnittstelle hat die alten Abfragen von InSIM nicht bereinigt``,
sagte Luca nach einer Stunde. Seine Stimme klang rauer als am Vorabend,
belegt vom feinen Staub der Nacht. „In Mailand laufen die Simulationen
ohne Schleifenüberwachung weiter. Sie generieren Agenten, die keine
Eingaben mehr erwarten.``

Phillips sah nicht auf. Sein Blick blieb fest verdrahtet mit der Linie
der Messmasten am Horizont. Eine Ader an seiner rechten Schläfe
pulsierte langsam, mit der sturen Gleichmäßigkeit eines Zeittaktes.

„Man kann eine Überlagerung nicht ewig aufrechterhalten``, fügte Luca
hinzu. „Irgendjemand muss die Messung beenden.``

Phillips griff nach der Blechtasse neben seinem Stuhl. Der Rand der
Blechtasse klirrte leise an seinen Zähnen. Er trank zwei Schlucke des
lauwarmen Wassers, stellte das Gefäß exakt an die Stelle zurück, an der
der feuchte Abdruck im Holz gezeichnet war, und strich sich mit der
Handfläche über die Stirn.

„Die Phänomenologie der Auslöschung``, sagte der alte Mann. Es war kein
Gesprächsbeitrag, eher das laute Aufrufen einer archivierten Fußnote.
„Wir haben in Rom Monate über die Invarianz des Schweigens diskutiert.
Aber die Gregoriana ist ein Raum aus Stein und Akustik. Dort hallt
selbst der Verzicht nach.``

Er stand auf, ohne das Holz des Stuhls knarren zu lassen, und ging in
den Schatten des Innenraums.

Gegen Mittag gingen sie nach Osten. Es gab keinen Pfad. Der Boden
bestand aus festem Sediment, unterbrochen von flachen Geröllfeldern,
deren Kiesel von einer feinen Ascheschicht überzogen schienen -- das
Überbleibsel alter vulkanischer Eruptionen, konserviert unter der
Wüstensonne.

Phillips ging mit gleichbleibender, langer Schrittlänge, die Arme lose
an den Seiten pendelnd. Er wich den Steinen nicht aus, er trat auf sie.
Das Knirschen des Gesteins unter den Sohlen war das einzige Geräusch in
der Senke.

An einer Senke aus hellerem Ton blieb der alte Mann stehen. Er bückte
sich nicht mit der Zögerlichkeit des Alters, sondern mit einer knappen,
funktionalen Beugung der Hüfte. Seine Finger griffen nach einem
faustgroßen, grauen Basaltstück. Die Oberfläche des Steins war rau,
übersät mit winzigen Poren, in denen sich der helle Sand festgesetzt
hatte.

Er hielt das Fragment auf der flachen Handfläche, ohne es zu drehen oder
zu prüfen. Er hielt es wie ein Messergebnis, das keiner Interpretation
bedurfte.

„Teilhard sah im Gestein die erste Stufe der Verdichtung``, sagte Luca,
der zwei Schritte hinter ihm verharrte. „Das Anreichern von Masse vor
dem Bewusstsein.``

Phillips hielt die Hand ruhig. Die Sonne warf den harten Schatten der
Finger auf den Stein.

„Die Steine sind keine Stufe``, sagte Phillips leise. Er ließ den Basalt
fallen. Der Aufprall auf den trockenen Ton war stumm, ein kurzer,
schwerer Druckimpuls im Staub. „Sie sind der Zustand nach dem Kollaps.
Sie bezeugen nichts, weil es für sie keine Alternativen gibt. Keine
Verzweigung. Keine Überlagerung.``

Er ging weiter, ohne sich umzusehen.

Die Tage verloren ihre Benennung. Sie unterteilten sich in das Wandern
der Schattensegmente auf den Dielen der Veranda, in die thermische
Ausdehnung des Wellblechs am Dach und in die kurzen Mahlzeiten aus
Hartbrot und Schafskäse, die sie schweigend auf der Brüstung einnahmen.

Am vierten Tag hatte sich die feine Staubschicht tief in die Poren von
Lucas Händen gesetzt. Die Haut an den Knöcheln war rissig geworden,
spröde wie altes Papier.

Er beobachtete die Flugbahn zweier Militärjets, die in großer Höhe durch
die Stratosphäre zogen. Die Maschinen selbst waren unsichtbar, nur zwei
messerscharfe, weiße Kondensstreifen schnitten den tiefblauen Himmel in
exakte Parallelen. Kein Schall erreichte den Boden. Die Triebwerke
blieben in der Höhe stumm.

„Martina hat sich aus dem Register des Parks abgemeldet``, sagte Luca in
die Stille hinein. Er saß auf dem warmen Holz und sah zu, wie die
Kondensstreifen sich langsam auflösten, breiter wurden und als grauer
Schleier vergingen. „Ihre Mutter sagte, sie sei nach Mailand gefahren.
Zu InSIM.``

Phillips bewegte die Finger seiner rechten Hand, als spiele er einen
stummen Akkord auf einer unsichtbaren Klaviatur.

„Sie suchen alle nach der Schließung des Modells``, antwortete der alte
Jesuit nach einer langen Pause. Seine Augen folgten den Rändern der
verblassten Streifen am Himmel. „Aber das ARS-System schließt nicht. Es
baut nur die Nebenschauplätze aus. Martina weiß das. Sie hat die
Grabungsberichte aus der Campania gelesen. Sie kennt den Unterschied
zwischen Konservierung und Erstarren.``

„Und Sie?{\kern0pt}``, fragte Luca. „Warum sind Sie hierher
gegangen?{\kern0pt}``

Phillips wandte den Kopf nicht. Die Haut seines Halses lag in tiefen,
trockenen Falten über dem Kragen.

„Man kann nicht gleichzeitig im Versuchsfeld stehen und die Messwerte
notieren, Luca``, sagte er, und seine Stimme verlor jede akademische
Tönung, wurde flach wie das verbleichende Licht über dem Becken. „Ich
habe jahrelang Dialoge berechnet. Ich habe Algorithmen beigebracht, wie
das Ungesagte strukturiert werden muss. Hier draußen gibt es keine
Struktur mehr. Nur das Signalrauschen des Hintergrunds.``

Er stand auf und trat an die Kante des Decks, die Hände hinter dem
Rücken verschränkt.

„Morgen``, sagte Luca leise zum Schatten des alten Mannes. „Morgen
breche ich die Reihe ab.``

Phillips antwortete nicht. Er stand an der Schwelle zum Staub, während
der Horizont die rote Farbe verlor und in das monotone Violett der
Wüstennacht überging.

\section{Kapitel 5: Die erste Frage}\label{kapitel-5-die-erste-frage}

Am siebten Tag zerbrach das Schweigen nicht in einem langen Gespräch,
sondern in kurzen, unvollständigen Sätzen, die sofort vom Wind verweht
wurden.

Die Sonne stand um elf Uhr grell über dem Hügelkamm. Die Luft über den
Basaltblöcken flimmerte so stark, dass die Konturen der alten Messmasten
am Horizont zu weichen, grauen Linien zu verschwimmen schienen. Es sah
aus wie ein Rechenfehler auf dem Monitor eines Oszilloskops.

Luca saß auf der untersten Stufe der Veranda, die Knie an den Brustkorb
gezogen. Er trug das Notizbuch auf dem Oberschenkel, aber der Stift lag
unangetastet in der Falte der Bindung. Die Dielen unter seinen Füßen
hatten die Wärme des Vormittags gespeichert und gaben sie in
gleichmäßigen, trockenen Wellen nach oben ab.

Phillips saß drei Stufen über ihm. Seine Hände lagen flach auf den
Oberschenkeln, unverändert, die Adern blau und hart abgezeichnet.

„Warum haben Sie die Gregoriana verlassen?", fragte Luca. Er sah nicht
nach oben. Sein Blick hing an einem gebleichten Echsenschädel, der halb
im roten Sand neben der Treppe vergraben lag.

Phillips bewegte den Kopf nicht. Die Haut an seinem Hals war straff
gezogen, als er den Atem anhielt.

„Weil die Dialoge korrupt wurden", sagte er schließlich. Die Stimme
klang kühler als an den Tagen zuvor, frei von Rauheit, fast mechanisch
artikuliert. „Wenn man einem System beibringt, Wahrscheinlichkeiten für
die nächsten Worte zu berechnen, verlernt man das Schweigen. Man
produziert nur noch Redundanz."

„Das InSIM-Labor braucht die Parameter für die Softwareagenten", sagte
Luca. „Mertens will die Prototypen in Pompeji einsetzen. Eine Simulation
von Bewusstsein, das nicht weiß, dass es in einer Schleife feststeht."

Phillips schlussfolgerte nicht. Er blickte über Lucas Kopf hinweg in das
grelle Weiß der Senke.

„Teilhard dachte, der Omega-Punkt sei das Ziel der Evolution", sagte der
alte Jesuit leise. „Die ultimative Verdichtung. Ein Zustand, in dem die
Materie so dicht gepackt ist, dass sie zu reinem Geist wird. Aber das
war die Theologie des zwanzigsten Jahrhunderts. Ein Irrtum der
Geometrie."

„Und was ist die Realität?", fragte Luca und drehte den unbenutzten
Stift in den Fingern.

„Das ARS-System rechnet nicht auf eine Verdichtung hin", antwortete
Phillips. Seine Finger krampften sich ganz leicht um das Gewebe seiner
Hose, dann entspannten sie sich wieder. „Es dehnt sich aus. Jede
Entscheidung erzeugt ein neues Feld. Der Omega-Punkt ist keine
Singularität, Luca. Er ist der Zustand, in dem keine Alternative mehr
berechnet werden muss. Nicht weil eine Welt übrig bleibt. Sondern weil
alle Fragen, die noch eine Berechnung verlangen, verschwunden sind."

Er deutete mit einer minimalen Bewegung des Kinns auf die Geröllhalde
vor der Veranda.

„Das sind die Steine", sagte er. „Kein Denkmal. Nur der Rest, wenn jede
Berechnung eingestellt wird."

Luca sah auf die grauen Blöcke aus erstarrtem Basalt. Die Hitze
spiegelte sich auf den glatten Abbruchkanten wie auf zerbrochenem Glas.

„Sie sprechen von einer Sackgasse", sagte Luca.

„Ich spreche von der Messung", entgegnete Phillips. „Die
Quantenschnittstelle kollabiert nicht von selbst. Man muss aufhören,
hinzusehen. Man muss das Labor verlassen."

„Und Martina?", fragte Luca. Die Frage fiel plötzlich in den Raum,
unvorbereitet, wie ein Gegenstand, der vom Tisch gleitet. „Sie ist nicht
nach Mailand gefahren."

Phillips reagierte nicht sofort. Der Schatten des Verandadachs schob
sich langsam über seine Stirn und tauchte seine Augenhöhlen in tiefe
Dunkelheit.

„Martina kennt den Unterschied zwischen einer Ruine und einer
Simulation", sagte er leise. „Sie sucht nach den Schichten unter der
Asche. Aber dort unten gibt es keine Welten mehr. Nur das Sediment."

Er stand auf. Die Gelenke seiner Knie gaben ein trockenes Knacken von
sich, das sofort im Flimmern der Luft verflog. Er ging nicht in die
Hütte, um Nahrung zu holen. Er trat an das Geländer, stützte sich mit
beiden Händen ab und sah zu, wie die Sonne den Horizont in ein hartes,
metallisches Bleichgelb tauchte.

„Du hast sieben Tage geschwiegen, Luca", sagte er, ohne sich umzudrehen.
„Aber du hast nicht zugehört. Du hast nur auf den Moment gewartet, an
dem du deine Fragen stellen kannst."

„Ich brauche eine Antwort für die Arbeit", sagte Luca.

„Es gibt keine Arbeit", antwortete Phillips. „Es gibt nur den
Versuchsaufbau. Und wir sind die Variablen, die man vergessen hat zu
löschen."

Er wandte sich ab und schritt durch die Fliegengittertür in den
schummrigen Innenraum.

Am achten Tag veränderte sich das Licht.

Ein trockener Nordwind war in der Nacht aufgekommen und hatte feinen,
rötlichen Staub aus den höheren Senken herangetragen. Der Himmel war
nicht mehr blau, sondern von einem fahlen, opaken Gelb, das die Sonne zu
einer matten Scheibe verblasste. Die Hitze war nicht mehr spürbar als
Strahlung, sondern als trockene Hülle, die auf der Haut lag wie ein
Panzer.

Sie saßen nebeneinander auf den Stufen. Keiner von beiden hatte das
Frühstück berührt. Die Plastikschale mit den getrockneten Datteln stand
unberührt zwischen ihren Füßen, überzogen von einem dünnen Schleier aus
feinstem Sand.

Luca hielt das Farbfoto des jungen Phillips in der Hand. Die Ränder des
Papiers hatten sich in der extremen Trockenheit der Wüste leicht nach
oben gewölbt.

„Er sieht aus, als wüsste er, was kommt", sagte Luca und reichte das
Bild hinüber.

Phillips nahm das Foto nicht. Er warf nur einen kurzen, flüchtigen Blick
auf den Papierabzug, als handele es sich um das Datenblatt eines längst
ausrangierten Geräts.

„Das war vor der Simulation", sagte der alte Mann. Seine Stimme war kaum
mehr als ein kühles Rauschen im Wind. „Bevor wir die Dialoggrammatiken
auf die Server gespielt haben. Damals glaubten wir noch, dass Sprache
ein Mittel sei, um Wirklichkeit zu teilen."

„Und heute?"

„Heute wissen wir, dass Sprache nur der Versuch ist, das Rauschen des
Systems zu überdecken. Jedes Wort ist eine Reduktion, ein schlechter
Kollaps."

Er griff nach seiner Wassertasse, trank nicht, sondern ließ die
Flüssigkeit langsam über die Kante auf die oberste Stufe laufen. Das
feuchte Holz färbte sich augenblicklich dunkel, und bereits nach wenigen
Sekunden begann die Stelle wieder zu verblassen, während der Staub das
Wasser aufsaugte.

„Du willst wissen, was du tun sollst, Luca?", fragte Phillips und sah
ihn seitlich an. Seine Augen wirkten in dem gelblichen Licht fast
transparent.

„Ja", sagte Luca.

„Nichts", antwortete der alte Jesuit. Er wandte den Blick wieder ab zur
staubigen Senke. „Setz dich hin. Beobachte das Abkühlen. Das ist die
einzige Bewegung, die keine neuen Welten erzeugt."

Sie saßen auf den Stufen, während der Staubwind in langen, dichten
Fahnen über das Sediment strich und die Spuren des Wagens auf der Piste
langsam, Satz für Satz, wieder auffüllte.

Am neunten Tag geschah das, was Luca später das zweite Licht nannte.

Es war nicht die Sonne. Es war auch nicht die matte Scheibe hinter dem
Staub. Es war ein Flirren am unteren Rand des Trümmerfelds, dort, wo die
Basaltblöcke in den hellen Ton übergingen. Luca sah es zuerst für eine
optische Täuschung an, eine Fata Morgana, die Hitze, die sich an den
Kanten der Steine brach. Aber das Flirren bewegte sich. Es kam näher. Es
nahm Form an.

Es war ein Mann.

Er ging nicht wie ein Mensch über den unebenen Boden. Er ging, als sei
der Boden eben. Seine Schritte setzten nicht auf, sie glitten. Er trug
eine dunkle, taillierte Jacke aus schwerer Baumwolle, denselben Schnitt,
den Phillips auf den Aufnahmen des Konsortiums von 1994 getragen hatte,
als die ersten Algorithmen für die Archäometrie von Sektor V
programmiert worden waren.

Er blieb drei Meter vor den Holzstufen stehen. Sein Gesicht war
unbeschädigt von der UV-Strahlung der Senke. Keine tiefen Kerben an den
Schläfen, keine Hornhautmembran über den Pupillen. Es war die
Physiognomie des Jesuiten aus dem Zeitungsausschnitt im Archiv von
Neapel: das schmale, akademische Kinn, die steile Stirnfalte eines
Mannes, der stundenlang über Matrizen brütete.

Der alte Phillips rührte sich nicht. Er saß auf der Veranda, die Hände
flach auf den Oberschenkeln, und sah den jüngeren Mann nicht an. Er sah
durch ihn hindurch, als sei er eine Störung im Signal, die man
ignorieren musste, um die Messung nicht zu verfälschen.

„Die Wahrscheinlichkeitsdichte in diesem Sektor ist extrem niedrig",
sagte der jüngere Phillips. Seine Stimme hatte keinen Nachhall, als
würde der Schall direkt von der Oberfläche der Bimssteinblöcke
absorbiert. „Ihr solltet den Brenner ausschalten. Das Abgas verändert
die Isotopenwerte im Sediment."

Luca richtete sich nicht auf. Seine Finger umfassten die raue, vom Harz
klebrige Kante des Geländers.

„Wir haben dich in den Außenbezirken gesucht", sagte Luca.

„Dort gibt es keine Überlagerungen mehr", sagte der jüngere Phillips. Er
blickte nicht zu Luca, sondern über das Wellblechdach hinweg auf die
Linie der Hügelkette. „InSIM hat die Abtastfrequenz gesenkt. Was ihr
Außenbezirke nennt, ist nur noch der Pufferbereich zwischen zwei
Iterationen. Wenn eine Variable ihre Rückmeldung verweigert, zieht das
ARS den Parameterbereich zusammen. Martina ist nicht gegangen. Sie hat
nur die Kohärenz mit dieser Speichereinheit verloren."

„Das ist ein theoretisches Konstrukt", sagte Luca.

„Es ist die Architektur, in der du sitzt", antwortete der jüngere
Phillips. Er hob die Hand, eine kurze, fast beiläufige Bewegung. Die
Haut an seinen Knöcheln war glatt, ohne Flecken. „Phillips hat in den
siebziger Jahren in Rom über die Logik der Verschränkung publiziert. Er
wusste, dass man ein Bewusstsein nicht simulieren kann, ohne ihm den
Austritt aus dem System zu erlauben. Wenn der Beobachter die Messung
nicht mehr bestätigt, kollabiert der Zustand nicht. Er verweilt in der
Überlagerung."

Der alte Phillips hob langsam den Kopf. Es war das erste Mal seit Tagen,
dass er eine Bewegung machte, die nicht dem Wind oder der Schwerkraft
folgte. Er sah den jüngeren Mann an - nicht als Vater den Sohn, nicht
als Forscher sein früheres Ich. Er sah ihn als das, was er war: eine
andere Perspektive auf denselben Grenzwert.

„Du bist nicht hier", sagte der alte Phillips.

„Ich bin auch hier", antwortete der jüngere Phillips. „Das ist der
Unterschied."

Luca sah zwischen beiden hin und her. Die Luft über dem Trümmerfeld
flimmerte stärker. Die Konturen der Basaltblöcke verschoben sich, als
würde die Landschaft selbst entscheiden, welche der beiden Gestalten sie
tragen wollte.

„Du willst, dass er die Messung abbricht", sagte Luca zum jüngeren
Phillips. „Du willst, dass er aufhört, hinzusehen."

„Nein", sagte der jüngere Phillips. „Ich will, dass er beides sieht. Das
Ende und die Fülle. Die Steine und die Möglichkeiten. Er hat sich
entschieden, nur eines zu sehen. Das ist sein Fehler. Nicht die Wüste.
Die Entscheidung."

Der alte Phillips stand auf. Die Bewegung war langsam, mechanisch, ohne
die Hast alter Menschen. Er trat an die Kante der Veranda, die Hände
hinter dem Rücken verschränkt, und sah hinaus in das Trümmerfeld.

„Teilhard", sagte der alte Phillips, „hat den Omegapunkt als Ziel
gedacht. Als Verdichtung. Als Vollendung. Er war Augustiner. Er war
Thomist. Er konnte nicht anders, als die Zeit auf ein Ende hin zu
denken. Das war seine Größe. Und sein Irrtum."

„Und was ist die Wahrheit?", fragte der jüngere Phillips.

„Es gibt nicht die Wahrheit", sagte der alte Phillips. „Es gibt zwei
Perspektiven auf denselben Grenzwert. Die eine sieht das Ende. Die
andere sieht die Fülle. Beide sind richtig. Keine ist ganz."

Luca trat einen Schritt vor. Die Hitze des Sediments drang durch die
Sohlen seiner Schuhe.

„Und wie entscheidet man, welche gilt?", fragte er.

„Man entscheidet nicht", sagte der alte Phillips. „Man ist die eine oder
die andere. Je nachdem, von wo man kommt. Je nachdem, was man verloren
hat. Je nachdem, was man noch sucht."

Der jüngere Phillips trat einen Schritt zurück. Mit diesem Schritt
schien der Kontrast seiner Gestalt zu schwinden, als verliere die
Auflösung seiner Konturen an Tiefe. Die dunkle Baumwolle seiner Jacke
floss in das Grau der Bimssteine über.

„Wenn du weiter nach ihr suchst", sagte die Stimme, die nun aus der
Richtung der Felszunge zu kommen schien, „verstärkst du nur das
Rauschen. Das ARS verarbeitet keine Trauer. Es verarbeitet nur
Koordinaten."

Dann war das Sediment leer. Kein Fußabdruck, kein verlagertes
Kieselgeröll. Nur das flache, monotone Flimmern der Mittagshitze auf dem
Basalt.

-\/-\/-

Der alte Phillips setzte sich wieder auf die oberste Stufe. Luca blieb
stehen.

„Er ist weg", sagte Luca.

„Er ist nie weg", sagte Phillips. „Er ist die andere Perspektive. Er
wird da sein, solange ich nur eine sehen kann."

„Dann sehen Sie beide?"

Phillips schwieg. Eine Minute. Zwei. Der Wind strich über die Dächer,
trug feinen Staub in die Ritzen des Holzes.

„Ich sehe das Ende", sagte er schließlich. „Ich kann die Fülle nicht
mehr sehen. Nicht weil sie nicht da wäre. Weil ich sie verloren habe. In
den Jahren. In den Entscheidungen. In der Arbeit."

„Und was hat Martina gesehen?"

Phillips sah auf seine Hände. Die Adern traten hervor wie blaue Kabel
unter der dünnen, pergamentartigen Haut.

„Martina hat beide gesehen", sagte er. „Deshalb ist sie gegangen. Nicht
in den Tod. Nicht in die Simulation. In die Möglichkeit. Sie hat die
Messung verweigert, weil sie wusste, dass jede Antwort eine Welt löscht.
Auch ihre eigene. Auch meine."

Luca setzte sich neben ihn auf die Stufe. Die Hitze des Holzes drang
durch den Stoff seiner Hose.

„Was ist dann der Omegapunkt?", fragte er. „Wenn er nicht Verdichtung
ist. Und nicht bloß Stille. Was ist er dann?"

Phillips sah in das Trümmerfeld hinaus. Die Basaltblöcke lagen da wie
eine Anordnung von Datenpunkten auf einem dunklen Feld, stumm, absolut,
unbeobachtet von den Kameras der Welt.

„Er ist der Punkt, an dem beide Perspektiven gleichzeitig gelten", sagte
der alte Mann. „Ohne dass sie sich versöhnen. Ohne dass sie sich
ausschließen. Er ist der Grenzwert der Perspektiven selbst. Nicht das
Ziel. Nicht das Ende. Der Rand, an dem man aufhört, sich zu entscheiden.
Und das Aufhören -- das ist das, was Teilhard nicht denken konnte. Weil
er Augustiner war. Weil er die Richtung für die Wahrheit hielt."

Luca sah ihn an. „Und Sie? Sind Sie auch Augustiner?"

Phillips schwieg. Ein Lächeln, kaum sichtbar, spielte um seine Lippen.
Es war kein freundliches Lächeln. Es war das Lächeln eines Mannes, der
erkannt hat, dass seine Frage falsch gestellt war.

„Ich bin Jesuit", sagte er. „Wir haben gelernt, in beiden Richtungen zu
denken. Aber das Lehren ist das eine. Das Sehen ist das andere."

Er stand auf. Die Gelenke gaben das trockene Knacken von sich, das Luca
inzwischen kannte. Er trat an das Geländer und sah hinaus in die Senke.

„Morgen", sagte Luca leise zum Schatten des alten Mannes. „Morgen breche
ich die Reihe ab."

Phillips antwortete nicht. Er stand an der Schwelle zum Staub, während
der Horizont die rote Farbe verlor und in das monotone Violett der
Wüstennacht überging.

In der Nacht geschah nichts. Kein Traum. Keine Erscheinung. Kein zweites
Licht. Nur die Stille, die über der Senke lag wie ein Bleidach. Und in
dieser Stille hörte Luca, zum ersten Mal seit Tagen, die zwei Sätze, die
Phillips gesagt hatte, noch einmal. Nicht als Worte. Als Zustände.

Das sind die Steine. Kein Denkmal. Nur der Rest, wenn jede Berechnung
eingestellt wird.

Er ist der Rand, an dem man aufhört, sich zu entscheiden.

Zwei Sätze. Der eine von DeLillo. Der andere von Van Dormael. Beide aus
demselben Mund. Beide wahr. Beide unvereinbar.

Und Luca begriff, dass er, wenn er die Wüste verließ, etwas mitnehmen
würde, das nicht in sein Notizbuch passte. Nicht eine Antwort. Sondern
die Frage, die beide Perspektiven offen ließ. Und das war, in Kapitel 5
der Zeitlosigkeit der Steine, die erste Frage, die nicht mehr gestellt
werden musste, weil sie mit der Wüste selbst identisch geworden war.

\section{\texorpdfstring{\hfill\break
}{ }}\label{section-1}

\section{\texorpdfstring{\textbf{Kapitel 6: Die
Steine}}{Kapitel 6: Die Steine}}\label{kapitel-6-die-steine}

Am neunten Tag trug der alte Mann feste Schuhe aus schwerem, rissigem
Leder und eine gefütterte Arbeitsjacke, die in der Hitze vollkommen
anachronistisch wirkte.

Als Luca um kurz nach sieben aus dem Türrahmen trat, stand Phillips
bereits dreißig Meter östlich der Veranda im Sand, die Hände in den
Jackentaschen versenkt, den Blick straff auf die Bruchkante der Senke
gerichtet. Er bewegte sich nicht, um auf Luca zu warten; er existierte
einfach an einem Punkt außerhalb der Hütte, als hätte ihn eine
Verschiebung im Geländeprofil dort abgelegt.

Luca griff nach den Stiefeln, die neben der Schwelle standen. Das Leder
war spröde, überzogen von feinstem Basaltstaub. Er schnürte sie, ohne
die Strümpfe aufzufalten.

Sie gingen nach Osten. Nicht entlang der verwaschenen Fahrspur, die
zurück zur Hauptstraße führte, sondern quer durch das freie Sediment.

Der Boden war kein Sandstrand. Es war ein erstarrtes Feld aus
zerkleinertem, porösem Vulkanit, durchsetzt mit scharfen
Quarzitsplittern, die unter den Sohlen ein trockenes, rhythmisches
Knirschen erzeugten. Die Luft schmeckte nach Schwefel und abgekühltem
Eisen -- derselbe Geruch, den Luca aus den klimatisierten
Untergeschossen des InSIM-Labors in Neapel kannte, wenn die
Hochleistungskühlungen der Quanten-Server nach zwanzig Stunden
Dauerbetrieb abgeschaltet wurden.

Phillips ging mit gleichmäßigem, mechanischem Schritt. Er wich den
größeren Felstrümmern nicht aus, sondern stieg mit starrer Hüfte darüber
hinweg. Seine Silhouette zeichnete sich scharf gegen das flache,
gelbliche Licht des Vormittags ab, dünn wie eine Schnittlinie auf einem
Diagramm.

Nach zweieinhalb Stunden veränderte sich die Geometrie der Landschaft.

Die flache Senke verengte sich zu einem breiten, baumlosen Kessel. Hier
gab es keine Vegetation mehr, keine verdorrten Dornensträucher, keine
Echsen im Schatten der Spalten. Vor ihnen lag eine Ansammlung von
Basaltblöcken -- hunderte, vielleicht tausende matter, dunkler Körper,
die flach auf der hellen Staubschicht lagen. Es sah nicht aus wie eine
natürliche Verwitterung. Es glich den Überresten einer städtischen
Struktur, deren Mauern und Dächer in einem einzigen Augenblick molekular
zertrümmert und flach auf den Boden gedrückt worden waren.

Phillips blieb am Rand des Kessels stehen. Er nahm die Hände aus den
Taschen und ließ sie schlaff an den Seiten hängen.

„Das ist kein Friedhof``, sagte er. Seine Stimme war leise, aber in der
absoluten Stille des Kessels klang jedes Wort wie ein sauber gesetzter
Impuls. „Ein Friedhof bewahrt Erinnerung. Das hier ist die Abwesenheit
von Signal.``

Luca trat an seine Seite. Die Hitze stieg aus dem dunklen Gestein auf,
unruhig flimmernd, so dass die weit entfernten Ränder des Kessels zu
schwanken schienen.

„Wir haben in Pompeji nach organischer Masse gesucht``, sagte Luca.
Seine Kehle war rau vom Staub. „Nach Abdrücken im Hohlraum der Asche.
Aber die Sensoren des ARS haben nur die Dichte der Steine gemessen. Das
System hat die Toten nicht rekonstruiert, Phillips. Es hat sie durch
mathematische Körper ersetzt.``

Phillips blickte nicht zu ihm. Seine Augenhöhlen lagen tief im Schatten
der Hutkrempe.

„Wenn man ein Ereignis bis ins kleinste Detail simuliert``, sagte der
alte Jesuit, „braucht man irgendwann keine Zeugen mehr. Die Messung
frisst das Subjekt. Was übrig bleibt, ist der Stein. Ein Körper, der
keine Fragen stellt und keine Parameter verändert.``

Luca bückte sich. Er griff nach einem faustgroßen Stück Basalt, das halb
im Staub eingebettet lag. Die Oberfläche war rau, porenreich, mit
winzigen grünlichen Olivinkristallen durchsetzt. Das Gestein fühlte sich
schwer an, überraschend kalt an der Unterseite, die dem Sand aufgeliegen
hatte, während die Oberseite glühte.

Er drehte den Stein in den Fingern, tastete die scharfen Kanten ab.

„Martina hat gesagt, dass die Steine in der Wüste Geduld haben``, sagte
Luca.

„Martina hat in Metaphern gedacht``, antwortete Phillips kaskadenartig,
ohne Satzmelodie. „Das war ihr Fehler. Die Steine haben keine Geduld.
Geduld setzt Zeit voraus. Steine sind der Zustand, in dem die Zeit
aufgehört hat, eine kontinuierliche Variable zu sein. Sie sind der
Endwert der Gleichung.``

„Der Omega-Punkt``, sagte Luca und ließ den Stein zurück in den Staub
fallen. Der Aufprall erzeugte einen dumpfen, klingenlosen Ton.

„Teilhard verwechselte den Punkt mit Gott``, sagte Phillips. Er machte
einen halben Schritt nach vorn, ohne in den Kessel hinabzusteigen. „Er
glaubte an eine geistige Vollendung. Aber der Omega-Punkt des
ARS-Projekts ist rein physikalisch. Es ist der Moment, in dem die
Rechenkapazität des Universums ausgeschöpft ist. Wenn alle Welten
berechnet sind und keine Wahrscheinlichkeit mehr übrig bleibt, um ein
weiteres Wort zu erzeugen. Dann erstarrt das System.``

„Und wir?{\kern0pt}``, fragte Luca.

Phillips sah ihn zum ersten Mal an. Seine Lider waren schwer, die
Pupillen zu Stecknadelkopfgröße verengt.

„Wir sind das Rauschen, das das System ausfiltert, bevor es sich
abschaltet. Wir sind die Zwischenwerte, die Asche.``

Er wandte sich nicht ab, um die Antwort abzuwarten. Er setzte sich
direkt auf einen der flachen Basaltblöcke am Kesselrand, zog die Knie
leicht an und legte die Ellbogen auf den Stoff seiner Arbeitsjacke.

Luca setzte sich drei Meter entfernt auf einen ähnlichen Stein. Die
Wärme des Vulkanits drang langsam durch den dichten Stoff seiner Hose.

Sie sprachen nicht mehr.

Die Sonne stieg höher, bis der Kessel keinen Schatten mehr warf. Die
Steine lagen da wie eine Anordnung von Datenpunkten auf einem dunklen
Feld, stumm, absolut, unbeobachtet von den Kameras der Welt.

Am späten Nachmittag gingen sie zurück. Der Rückweg dauerte länger. Die
Schatten der beiden Männer dehnten sich über den rötlichen Sand aus,
lange, feine Linien, die vor ihnen herliefen und die Bruchkanten der
Steine wie schwarze Tinte auffüllten.

Als sie die Hütte erreichten, war das Licht im Westen in ein trübes,
violetteres Sediment übergegangen.

Sie tranken lauwarmes Wasser aus den Blechbechern auf der Veranda.
Keiner von beiden machte Licht im Innenraum.

Am zehnten Tag veränderte sich die Textur des Schweigens.

Als Luca am Morgen hinausging, saß Phillips auf der obersten Stufe. Er
trug die Arbeitsjacke nicht mehr, sondern das einfache, ausgebleichte
Hemd der ersten Tage. In seiner rechten Hand hielt er einen ausgebauten
Sensor aus der ersten Versuchsreihe der Messmasten, dessen
Kabelanschluss sauber abgetrennt worden war.

Er drehte das Bauteil zwischen Daumen und Zeigefinger.

Luca setzte sich auf die unterste Stufe. Er stellte die Fragen nicht
mehr. Er sah zu, wie das Sonnenlicht sich in dem kleinen Glasfenster des
Sensors brach, ein winziger, hochintensiver Lichtpunkt, der über das
Holz der Stufe wanderte.

„Man kann das Experiment nicht abbrechen, Luca``, sagte Phillips leise.
Es klang nicht wie eine Belehrung, sondern wie die Feststellung einer
physikalischen Konstante.

„Ich weiß``, antwortete Luca.

„Das ARS rechnet weiter, ob wir in Neapel vor den Monitoren sitzen oder
hier draußen im Staub. Es erzeugt Milliarden von Welten, in denen
Martina nach Hause fährt, in denen sie in Mailand aus dem Zug steigt, in
denen sie nie existiert hat.``

Er ließ den Sensor in die Sandspur neben der Treppe fallen. Das Bauteil
verschwand halb in der roten Asche.

„Was tun wir dann hier?{\kern0pt}``, fragte Luca nach einer langen
Pause.

Phillips lehnte den Hinterkopf gegen die raue Holzwand der Hütte und
schloss die Augen.

„Wir bleiben in der einzigen Welt, in der wir nicht mehr antworten
müssen.``

Sie saßen da, während die Hitze des zehnten Tages über die Senke stieg
und die Wüste vor ihnen langsam zu einem einzigen, spiegelnden Feld aus
grellem Weiß zusammenschmolz.

\section{\texorpdfstring{\hfill\break
}{ }}\label{section-2}

\section{\texorpdfstring{\textbf{Kapitel 7: Die
Rekonstruktion}}{Kapitel 7: Die Rekonstruktion}}\label{kapitel-7-die-rekonstruktion}

Am vierzehnten Tag veränderte ein mechanisches Geräusch die Dichte der
Luft.

Es war kein plötzlicher Lärm, sondern eine tiefe, langsame Frequenz, die
weit im Süden aus dem Hitzeflimmern der Fahrspur wuchs. Ein alter Sedan,
dessen ausgebleichter Lack die Farbe von getrocknetem Lehmschlamm
angenommen hatte, schob eine flache Bugwelle aus grauem Staub vor sich
her. Der Motor lief unruhig, mit dem unregelmäßigen Taktieren eines
Ventils, das kurz vor dem Festfressen steht.

Phillips saß auf der zweiten Holzstufe der Veranda. Er bewegte den Kopf
nicht. Seine Hände lagen flach auf den Oberschenkeln, starr, als spürte
er die Vibrationen der Achsen bereits durch das vulkanische Fundament.

Der Wagen hielt fünf Meter vor den Stufen. Der Kühler zischte leise, ein
dünner Strahl Kühlflüssigkeit verdampfte stotternd auf dem heißen Blech
des Unterbodens.

Die Fahrertür schwang auf.

Eine Frau stieg aus. Sie war Mitte vierzig, trug eine dunkelblaue
Baumwolljacke über einem verwaschenen Shirt und feste Leinenschuhe. Ihr
Haar war streng nach hinten gebunden, von feinstem Sandstaub mattgrau
gepudert. Sie trug kein Gepäck, nur eine abgewetzte Reisetasche aus
braunem Segeltuch, deren Reißverschluss zur Hälfte offen stand. In dem
Spalt war das gequetschte Papier eines alten Fahrplans aus Neapel zu
erkennen.

Sie blieb an der Stoßstange stehen. Ihre Augen wanderten über die
verwitterten Bretter der Hütte, glitten stumm über Phillips hinweg und
blieben schließlich an Lucas Schultern hängen.

„Der Zug aus San Diego hatte vier Stunden Verspätung``, sagte sie.
Herkunft und Tonfall ihrer Stimme trugen die genaue Kadenz jener
akademischen Zirkel Neapels, in denen Phillips einst gelehrt hatte --
kühl, präzise, bar jeder dramatischen Intonation. „Die Verbindung in
Indio existiert nur noch auf den Netzplänen.``

Phillips antwortete nicht. Er fixierte einen Punkt weit hinter dem Heck
des Wagens, dort, wo die Fahrspur im Sediment auslief.

Sie sah ihn an, ohne den Kopf zu neigen. „Sie nimmt die Medikamente
nicht mehr, Michael. Der Arzt in der Klinik sagt, die Neuronenmuster
reagieren nicht auf die Korrektur. Sie hat aufgehört, nach dir zu
fragen.``

Sie griff nach ihrer Tasche, ging an den Stufen vorbei und trat durch
die offene Tür ins Innere der Hütte. Ihre Schritte auf den Holzdielen
klangen hart, ungepolstert. Ein kurzes, metallisches Klirren folgte, als
sie die Tasche auf den Tisch stellte -- genau auf jene Stelle, an der
bis zum Vortag die Aktenordner der Versuchsreihe gelegen hatten.

Als sie wieder ins Licht der Veranda trat, trug sie das Segeltuch nicht
mehr. Sie setzte sich auf die Stufe unterhalb von Phillips. Ihr Rücken
berührte fast seine Knie, doch keiner von beiden suchte den Kontakt.

Sie drehte den Kopf leicht zu Luca. Die Pupillen waren dunkel, von
feinen Rötungen durchzogen, die Augenlider scharf gezeichnet.

„Er hat mir erzählt, dass Sie hier sein würden``, sagte sie. „Ein
Assistent aus der Messtechnik. Jemand, der an die Rekonstruktion
glaubt.``

„Ich bin nicht wegen der Messungen hier``, sagte Luca.

„Warum dann?{\kern0pt}``, fragte sie. Sie strich mit zwei Fingern über
das Holz der Stufe, schabte eine feine Kruste aus getrocknetem Staub ab.
„Um zuzusehen, wie ein System ausbrennt? Um Protokoll zu
führen?{\kern0pt}``

„Er hat das Labor nicht verlassen, weil die Daten falsch waren``, sagte
Luca. Seine Kehle spannte sich an. „Das ARS hat nicht simuliert,
Martina. Es hat die Möglichkeiten abgeschafft. Jede Schaltung hat eine
Welt gelöscht.``

Ein kurzes, lautloses Lächeln schnitt über ihr Gesicht, kühler als die
Schatten unter dem Wellblechdach.

„Ich heiße nicht Martina``, sagte sie leise. „Martina ist die Variable
aus den Simulationen. Die Tochter, die in der Schicht von
siebenundsiebzig im Bus nach Rom saß. Diejenige, deren Parameter Sie in
den Rechenzentren von InSIM alle sechs Wochen korrigiert haben.``

Luca sah sie an. Der Wind fasste in das dünne Kleid, drückte den Stoff
gegen ihre Rippen.

„Sie existiert im Datensatz``, sagte er.

„Der Datensatz ist eine mathematische Leiche``, antwortete sie, ohne die
Stimme zu heben. „Ein Vektor aus Wahrscheinlichkeiten, konstruiert aus
seinen Vorlesungsskripten und den Beichtprotokollen der Jesuiten. Mein
Vater hat nie gelernt, zwischen einer Person und ihrer Wellenfunktion zu
unterscheiden. Für ihn war das Verschwinden nur eine Zustandsänderung,
die man durch genügend Iterationen rückgängig machen kann.``

Phillips rührte sich nicht. Das Licht brach sich in seinen verwaschenen
Hornhauträndern, starr und spiegelnd.

„Wir haben das Signal verloren``, sagte der alte Mann plötzlich. Es war
das erste Mal seit drei Tagen, dass seine Stimme den Raum füllte. Sie
klang flach, bar jeder Resonanz, wie eine Durchsage über einen kaputten
Lautsprecher. „Am vierundzwanzigsten August. Die Sensoren im Sektor Vier
zeigten nur noch Rauschen. Keine Masse. Kein Körper. Nur noch der
Zerfall der Kohärenz.``

Die Frau sah auf ihre Hände. Ihre Fingernägel waren kurz geschnitten,
die Knöchel vom Staub weißlich eingefärbt.

„Ich war am vierundzwanzigsten August in Pompeji, Michael``, sagte sie
ruhig. „Ich bin durch die Ausgrabungen gegangen. Ich bin an den
Gipsabgüssen vorbeigelaufen. Ich habe auf mein Telefon gesehen. Es gab
kein Signal, weil die Masten im Tal abgeschaltet waren. Nicht, weil die
Welt aufgehört hat zu existieren.``

Sie stand auf. Das Holz der Stufe ächzte unter der leichten Entlastung.

„Der Wagen hat keinen Treibstoff mehr für die Rückfahrt``, sagte sie zu
Luca. „Ich habe unterwegs in Ocotillo angehalten, aber die Zapfsäulen
sind gesperrt. Sie haben dort vor drei Wochen den Strom abgedreht.``

Sie ging nicht zurück in die Hütte. Sie spazierte langsamen Schrittes um
den verstaubten Sedan herum, strich mit der Handfläche über das heiße
Blech der Motorhaube und blieb am Heck stehen. Aus der Wüste stieg die
Hitze des späten Nachmittags auf, dicke, träge Schichten aus flimmernder
Luft, die das Bild der Bergkette im Norden in verzerrte, waagerechte
Streifen auflösten.

Sie saßen auf der Veranda, bis das rotviolette Licht der Abenddämmerung
in das graue Sediment überging. Kein Stern war zu sehen; die Atmosphäre
über der Senke war gesättigt von feinstem Schwebstaub.

Am nächsten Morgen war der Wagen weg.

Es gab kein Motorgeräusch. Keine Reifenspuren im Staub. Der Platz vor
der Veranda war leer, als hätte der Ton des alten Sedans nie die Luft
verdichtet. Nur die beiden tiefen Rillen im Vulkanit, in denen das
Kühlwasser am Vortag versickert war, zeichneten sich als dunkle,
verkrustete Linien im Sand ab.

Luca trat auf die Stufe. Der Wind stach ihm kalt in die Augen.

Phillips saß an derselben Stelle wie am Vorabend. Vor seinen Füßen lag
das kleine Glasfenster des abgetrennten Sensors. Die Morgensonne
spiegelte sich darin -- nicht als Lichtpunkt, sondern als stumpfe, matte
Fläche, über die eine Ameise träge hinwegkrabbelte.

„Sie wird nicht wiederkommen``, sagte Luca.

Phillips wandte den Kopf nicht. Seine Wangen wirkten eingefallen, die
Haut über den Jochbeinen gespannt wie Pergament.

„Die Wahrscheinlichkeit für eine Rückkehr liegt innerhalb der
Fehlertoleranz``, sagte er leise. „Wenn man ein System lange genug
beobachtet, werden die Ränder scharf. Es gibt keine alternativen
Zustände mehr. Keine Tochter. Keine Stadt unter der Asche.``

Er hob die Hand, langsam und mit einer beinahe mikroskopischen Bewegung,
und deutete auf das Feld aus dunklen Basaltblöcken, das sich jenseits
der Senke in der Morgensonne dehnte.

„Die Steine messen nicht mehr``, sagte er.

Luca setzte sich auf die Stufe neben ihn. Er zog die Knie an und legte
die Arme darüber. Die Kälte des Steins drang durch seine Sohlen. Sie
sprachen nicht mehr. Das Flimmern über dem Feld begann erneut.

\section{\texorpdfstring{\hfill\break
}{ }}\label{section-3}

\section{\texorpdfstring{\textbf{Kapitel 8: Die Stille der
Tochter}}{Kapitel 8: Die Stille der Tochter}}\label{kapitel-8-die-stille-der-tochter}

Am sechzehnten Tag hatte sich das Flimmern nicht aufgelöst.

Der Sedan stand noch immer am Fuß der Veranda, schief abgesackt im
Lockersand der linken Baugruppe. Die Hitze stieg in sichtbaren Wellen
von den Kotflügeln auf, wie das Entweichen von Druck aus einer
unvollständig isolierten Kammer.

Die Frau saß auf der untersten Holzstufe. Ihr Nacken war gebeugt, die
Sehnen unter der matten Haut straff gezogen wie Drähte in einem Gehäuse.
Sie sah nicht auf das basaltische Trümmerfeld hinaus, sondern auf den
Raum zwischen ihren eigenen Schuhen, wo sich feinstes Vulkanit-Sediment
in den Falten des Segeltuchs sammelte.

Luca trat aus der Türöffnung. Der Türrahmen roch nach verharzter Kiefer
und Altöl. Er setzte sich drei Stufen oberhalb ab, das Holz gab mit
einem trockenen Knacken nach.

„Sie haben die Zündung nicht unterbrochen``, sagte er.

Sie bewegte den Kopf nicht. „Das Ventil klemmt. Wenn der Druck abfällt,
zieht der Kopf Luft. Es ist kein thermisches Problem, sondern eine
mechanische Störung.``

Ihre Stimme besaß die charakteristische Trockenheit des neapolitanischen
Universitätsviertels, aber da war eine Schicht aus Staub darunter, als
hätte sie die Vokale durch ein Sieb gegossen.

Phillips erschien Minuten später. Er trug noch immer das helle Hemd mit
den dunklen Verfärbungen unter den Achseln, deren Ränder salzweiß
ausgeblüht waren. Er hielt sich nicht am Pfosten fest. Er glitt an der
Wand entlang, bis er saß, die Ellbogen auf den Kniegelenken verkeilt,
die Finger ineinandergefügt wie bei einer Übung zur Rekonstruktion von
Kristallgittern.

Das Schweigen war keine Absenz von Schall. Es war eine dichte,
physikalische Materie, gebildet aus dem kontinuierlichen Taktieren des
abkühlenden Blechs und dem Piezosumpf der Insekten im Trockengras.

„Die Kurve im Sektor Neun war nie geschlossen``, sagte Phillips nach
einer Stunde. Seine Pupillen verengten sich im grellen Reflex des
Sedimentbodens. „Teilhard hat das Omega-Punkt genannt. Aber es gab keine
Konvergenz. Nur die Zunahme der Entropie in den Randbereichen.``

Die Frau strich mit dem Zeigefinger über das Leder ihrer Sandale. Eine
Spur aus grauem Schmutz blieb an der Beuge hängen.

„In der Klinik in Capodichino gab es einen Raum mit zwei Fenstern``,
sagte sie, ohne aufzusehen. „Man konnte das Meer sehen und die
Einflugschneise des Flughafens. Sie hat den ganzen Nachmittag dort
gesessen und die Frequenzen der Turbinen gezählt. Sie hat gesagt, bei
jeder Landung wird eine Ausfallstraße in der Matrize gelöscht.``

Luca sah auf ihre Schulterblätter. „Das waren die Beruhigungsmittel.``

„Nein``, sagte sie kühl. „Das war die Syntax. Ihre Syntax. Er hat ihr
beigebracht, die Welt als einen Baum von Fallentscheidungen zu lesen.
Wenn man lange genug zählt, bleibt kein Stamm mehr übrig. Nur noch
Abzweigungen.``

Der siebzehnte Tag.

Die Sonne stand um elf Uhr als blasser, feuriger Punkt hinter einer Wand
aus atmosphärischem Schwebstaub. Der Horizont hatte seine Kontur
verloren; das Gestein und der Himmel verschmolzen zu einem einzigen,
bleichen Kontinuum.

Die Frau hatte ihre Jacke abgelegt. Das Shirt war an den
Schulterblättern durchgefeuchtet, die Konturen ihrer Wirbelsäule
zeichneten sich scharf ab wie die Segmente eines fossilen Rückgrats.

„Die Abgüsse in Pompeji sind hohl``, sagte sie zu dem Raum zwischen
ihren Kniekehlen. „Das Ministerium hat in den Sechzigern versucht,
Polyurethan einzuspritzen. Um die organischen Reste zu konservieren.
Aber das Gewebe war längst fort. Sie haben nur die Form des Hohlraums
versiegelt.``

Phillips rührte sich nicht. Ein feiner Schweißfilm glänzte auf seiner
Stirn, ohne abzuperlen.

„Der Hohlraum ist der Zustand vor der Messung``, sagte er. Die Stimme
klang wie trockenes Papier, das über Beton gezogen wird. „Das ARS hat
die Hohlräume kartiert. Mehr nicht. Wenn man Flüssigkeit hineinpumpt,
zerstört man die Überlagerung.``

„Sie haben Plastik hineingepumpt, Michael``, sagte sie leise. „Es stinkt
noch heute nach Lösungsmitteln, wenn die Sonne auf die Vitrinen fällt.``

Luca beugte sich vor. Er spürte die Hitze der Holzdielen durch die
Oberschenkel. „Warum sind Sie von der Küste weggegangen?{\kern0pt}``

Sie drehte das Gesicht leicht zu ihm. Die Rötung in ihren Augen war
tiefer geworden, feine Verästelungen aus geplatzten Kapillaren.

„Weil die Datenbanken voll waren``, sagte sie. „InSIM hat die
Rechenkapazität im Mai halbiert. Die Server in Rom verbrauchen zu viel
Kühlwasser. Sie haben die Nebenpfade gelöscht. Martina war ein
Nebenpfad.``

„Und Sie?{\kern0pt}``, fragte Luca.

Sie sah ihn an, starr, ohne zu blinzeln. „Ich bin der Restfehler. Die
Zahl, die man stehen lässt, weil das Runden mehr Rechenzeit kostet als
der Speicherplatz.``

Am Abend erfüllten die Grillen das Trockengestrüpp hinter der Hütte mit
einer Intensität, die den eigenen Herzschlag überlagerte. Sie aßen kalte
Bohnen aus der Dose, ohne Besteck, mit den Fingern die fettige Schicht
von den Blechrändern schabend.

Phillips nahm kein Essen an. Er saß im Schatten der Innenwand, das
Gesicht der dunklen Ecke zugewandt, in der einst die Terminals gestanden
hatten.

„Es gibt eine Formulierung bei Loyola``, murmelte er in die Dämmerung
hinein. „Die dritte Art der Demut. Wenn das Subjekt wünscht, als tot
betrachtet zu werden, damit die Ordnung der Dinge nicht durch seine
Gegenwart gestört wird.``

„Das ist kein Loyola, Vater``, sagte die Frau vom Türrahmen aus. Sie
hielt die leere Blechdose in beiden Händen, als wöge sie deren Masse ab.
„Das ist die Dokumentation der Rechenkerne. Seite vierzig. Korrektur der
Drift.``

Der achtzehnte Tag.

Am Morgen war der Wind umgeschlagen. Er kam nun direkt aus dem Becken
von Ocotillo, beladen mit dem Geruch von verbranntem Gummi und trockenem
Salz.

Der Sedan stand unbewegt da. Der Reifen vorne rechts hatte über Nacht
Luft verloren; die Felge ruhte jetzt direkt auf einem flachen
Basaltblock.

Die Frau stand auf der Freifläche vor der Veranda. Sie trug die braune
Segeltuchtasche unter dem Arm. Die Bänder hingen herunter und schleiften
durch den feinen Sedimentstaub.

„In der dritten Schicht des ARS gab es einen Parameter für das
Schweigen``, sagte Luca. Er stand auf der obersten Stufe, die Hand flach
gegen den ausgetrockneten Türpfosten gepresst. „Ein Schwellenwert. Wenn
die Entropie einen bestimmten Wert überschritt, stellte der Algorithmus
die Ausgabe ein.``

Sie drehte sich nicht um. „Das war kein Algorithmus. Das war die
Erschöpfung des Personals.``

Sie ging los.

Sie nutzte nicht die Fahrspur, die der Sedan im Staub hinterlassen
hatte, sondern schnitt geradlinig durch das Trümmerfeld aus dunklen
Basaltblöcken. Ihre Leinenschuhe setzten mit einem kurzen, trockenen
Knirschen auf den scharfkantigen Steinen auf.

Luca machte zwei Schritte nach vorn. Die Hitze schlug ihm wie eine feste
Platte gegen die Brust.

„Lass sie``, sagte Phillips hinter ihm.

Die Stimme des alten Mannes war so leise, dass sie fast im Zischen des
Windes unterging.

Luca sah nicht zurück. „Sie geht in die Senke. Dort gibt es keinen
Schatten. Keine Straßen.``

„Sie geht nicht in die Senke``, sagte Phillips. Er war nicht
aufgestanden. Er saß auf dem Dielenboden, die Hände um die Knie
geschlossen, den Kopf leicht geneigt, als lausche er dem Geräusch einer
entfernten Hydraulik. „Sie geht nur aus dem Blickfeld der Sensoren. Das
ist der Unterschied.``

Luca beobachtete die dunkle Silhouette der Frau. Sie wurde kleiner,
nicht fließend, sondern in Rucken, verzerrt durch die aufsteigenden
Schichten der Hitzekonvektion. An einem Punkt, etwa zweihundert Meter
entfernt, wo zwei mächtige Vulkanitblöcke ein schmales Tor bildeten,
löste sich ihre Kontur auf. Das dunkelblaue Baumwollgewebe ihrer Jacke
schmolz mit den dunklen Tönen des Gesteins zusammen.

Ein Flimmern, kurz und scharf, als hätte eine Linse die Brennweite
korrigiert.

Dann war da nur noch der Basalt.

Luca blieb auf der Stufe stehen. Er spürte, wie der feine Schwebstaub
sich auf seine Lippen legte, trocken, mit einem metallischen
Nachgeschmack von Schwefel und veralteter Elektronik.

„Sie war hier``, sagte er. Es war keine Frage.

Phillips antwortete nicht sofort. Er langte nach unten und hob das
kleine Glasfenster des abgetrennten Sensors auf. Er drehte es langsam
zwischen Daumen und Zeigefinger. Die Oberfläche war von Mikrokügelchen
aus geschmolzenem Quarz übersät.

„Das System hat keine Historie mehr``, sagte der alte Mann. Seine Augen
waren völlig starr, gerichtet auf die matten Lichtreflexe im Glas. „Wenn
man einen Zustand nicht mehr liest, hat er nie stattgefunden. Das ist
die Eleganz der Architektur.``

Luca setzte sich auf die heiße Stufe. Er legte die Handflächen auf die
Knie, genau in die Abdrücke, die Phillips vor drei Tagen hinterlassen
hatte.

Die Steine vor ihnen veränderten ihre Farbe nicht. Sie absorbierten das
Licht, stumm, schwer, ohne jede Neigung zur Resonanz.

\section{\texorpdfstring{\textbf{Kapitel 9: Die leere
Seite}}{Kapitel 9: Die leere Seite}}\label{kapitel-9-die-leere-seite}

Am neunzehnten Tag war die Hütte von einem trockenen, mehligen Knirschen
erfüllt.

Es war das Gebälk, das unter der Vertikalen der Mittagssonne nachgab.
Luca lag mit offenen Augen auf dem Gestell der Pritsche. Der Kissenbezug
roch nach altem Talg und dem unorganischen Staub der Lüftungsschächte
von InSIM. Er richtete sich auf, ohne die Füße auf den Boden zu setzen.
Das Licht schob sich in schmalen, messerscharfen Streifen durch die
Ritzen der Wellblechwand, Schnittlinien im Dämmerlicht der Kammer.

Er trat ins Freie. Die Luft hatte die Viskosität von abgekühltem Öl.

Phillips saß auf der Veranda, die Unterarme auf den Oberschenkeln
verschränkt, den Blick auf den blassen Schotterstreifen gerichtet, wo am
Vortag der Schatten des Sedans geendet hatte. Der Wagen stand noch da,
seine Karosserie tief im Sediment eingesunken, als habe das Eigengewicht
des Metalls den lockeren Boden unter dem Basalt vollkommen verdichtet.

„Sie ist nicht durch den Sektor Vier gegangen``, sagte Luca.

Phillips antwortete nicht. Seine Finger tasteten langsam die Naht seiner
Hose entlang, hin und zurück, mit dem gleichmäßigen Taktieren eines
Relais im Standby-Modus.

Luca ging um das Gebäude. Die Rückwand war roh beschlagen, die
Schnittkanten der Bretter vom Ausblühen des Harzes weißgrau verkrustet.
Er lief die flache Böschung hinab bis zu jener Felszunge, an der die
Versuchsreihe im Juni begonnen hatte. Der Staub nahm seine Schritte
lautlos auf. Keine Spur, kein Abdruck im lockeren Bims. Er rief ihren
Namen nicht. Die Frequenz einer menschlichen Stimme wäre in der Weite
des Beckens sofort verdampft, ohne Echo, ohne Widerhall gegen die
erstarrten Magmaformationen.

Im Inneren der Hütte war die kleine Kammer hinter der Küche verwaist.

Das Bettzeug lag strammgezogen über der Rosshaarmatratze. Kein Kissen
war eingedrückt, keine Falte verriet die Kontur eines ruhenden Körpers.
Auf der rohen Fichtenplatte des Tischchens lag nur ein querformatiges
Notizbuch aus italienischer Fertigung, der Einband aus abgewetztem,
schwarzem Wachstuch, die Ecken von der Luftfeuchtigkeit der Küste leicht
aufgeworfen.

Es lag aufgeschlagen da.

Luca trat an den Tisch. Seine Hand verharrte zwei Zentimeter über dem
Papier. Die rechte Seite war weiß, von einer fast aggressiven,
unberührten Dichte. Kein Graphitkrümel, keine Kerbe von einem zu fest
aufgedrückten Stift.

Er wandte sich um. Phillips stand im Türrahmen. Der Schatten des alten
Mannes war kurz, eine dunkle, verdichtete Masse auf den Dielen.

„Sie hat nichts hinterlassen``, sagte Luca.

Phillips trat näher. Sein Atem klang rasselnd, wie Sand, der durch ein
feines Messingsieb gleitet. „Es gibt im Modell von Everett keine
Hinterlassenschaften. Wenn ein Zweig der Wellenfunktion seine Kohärenz
verliert, nimmt er seine Ränder mit.``

Luca griff nach dem Buch und blätterte zurück. Die Blätter knisterten
trocken.

Auf den ersten Seiten befanden sich Tuschezeichnungen. Keine
Landschaften, sondern schematische Schnittzeichnungen: die
Schichtenfolgen der Ausgrabungsfelder von Pompeji, Grundrisse der Regio
V, schraffierte Quadrate mit Millimeterangaben zu Vulkanit-Sedimenten.
Dazwischen, in der feinen, steilen Handschrift der Frau, mathematische
Matrizen für den Zerfall von Isotopen.

Er erreichte die letzte beschriebene Seite.

Da war kein langer Text, keine Rechtfertigung. Nur eine einzige,
isolierte Zeile am oberen Rand, mit dünner Tinte geschrieben, die an den
Rändern der Buchstaben mikroskopisch ausgefranst war:

\emph{In einer anderen Iteration habe ich die Parameter nicht
bestätigt.}

Luca schloss das Buch. Das Wachstuch fühlte sich fettig an unter seinen
Fingerspitzen.

„Sie hat das System nicht verlassen``, sagte er.

Phillips sah auf die leere Tischplatte. Seine Augen wirkten trüb, als
überziehe eine dünne Membran aus Hornhaut das Licht seiner Pupillen.
„Sie hat die Messung verweigert. Das ist die einzige Form der Löschung,
die der Algorithmus akzeptiert. Wer den Zustand nicht rückmeldet, fällt
aus der Matrix. Wir haben das im Versuchsfeld Sechs beobachtet.
Dreihundert Sensoren. Sie zeigten nicht
\textquotesingle Null\textquotesingle{} an. Sie zeigten
\textquotesingle Kein Empfänger\textquotesingle.``

„Sie war eine Person, Michael. Kein Sensor.``

„Das ist ein semantischer Unterschied, Luca. Eine Frage der Skalierung.
Das ARS unterscheidet nicht zwischen einer biologischen Schnittstelle
und einer Diode, wenn die Datenströme abbrechen.``

Der zwanzigste Tag.

Die Mittagshitze drückte den Rauch des Kochers flach auf den Boden. Sie
hatten das Gestell vor die Veranda gestellt. Eine Blechdose mit Wasser
stand auf dem Brenner, das Wasser stand kurz vor dem Sieden, ohne Blasen
zu bilden, träge und stumm.

Luca sah auf die Steine. Das Trümmerfeld erstreckte sich bis zum
Horizont wie das erstarrte Abbruchfeld einer antiken Nekropole. Die
Basaltblöcke wirkten im blendenden Licht völlig flach, ohne Tiefe, wie
auf das Panorama der Wüste gemalt.

„In Pompeji``, sagte Luca, „gibt es die Casa del Fauno. In einem der
Hohlräume haben sie den Abdruck eines hölzernen Schrankes gefunden. Die
Maserung des Holzes war im Bimsstein konserviert. Das Holz selbst war
seit zweitausend Jahren Staub.``

Phillips rührte sich auf der Stufe nicht. Seine Hände lagen flach auf
seinen Knien, die Haut dünn wie Pergamentpapier über den bläulichen
Venenbögen.

„Das Material verfällt``, sagte er leise. „Was bleibt, ist die Geometrie
der Abwesenheit. Das ARS war nie eine Rekonstruktion des Lebens, Luca.
Es war die Kartierung der Hohlräume. Wir haben die Konturen der Toten
nachgezeichnet, bis das Signal stark genug war, um die Lebenden zu
ersetzen.``

Luca sah ihn an. „Und Martina?{\kern0pt}``

„Welche Martina?{\kern0pt}``, fragte der alte Mann. Seine Stimme klang
nicht zynisch. Sie war flach, völlig gereinigt von Emotion, wie die
maschinelle Sprachausgabe eines alten Terminals. „Die Tochter, die in
Rom Vorlesungen besuchte? Die Frau, die gestern durch das Trümmerfeld
ging? Oder die Variable im Algorithmus, die wir jeden Abend um sechs Uhr
abgeglichen haben? Sie sind alle im System gedämpft worden. Eine nach
der anderen.``

Phillips stand auf. Seine Bewegungen waren extrem langsam, als koste es
ihn eine unermessliche Anstrengung, die Masse seines Körpers gegen die
Schwerkraft der Senke aufzurichten.

Er ging die Holzstufen hinab. Seine Sohlen hinterließen keine tiefen
Abdrücke im heißen Staub. Er ging zwei Schrittlängen in das Trümmerfeld
hinein und blieb vor einem hüfthohen, scharfkantigen Vulkanitblock
stehen.

Er berührte den Stein nicht. Er neigte nur den Kopf, als lausche er
einer Frequenz, die knapp unterhalb des menschlichen Hörbereichs lag.

„Sie messen nicht mehr``, wiederholte der alte Mann. „Die Steine haben
die Eigenschaft verloren, Daten aufzunehmen. Wir sind am Nullpunkt,
Luca. Das Omega ist erreicht, aber es gibt keine Erhebung. Nur die
Stille des Materials.``

Luca folgte ihm nicht. Er blieb auf der Veranda sitzen, die Ellbogen auf
den Knien.

Die Luft über den Basaltblöcken zitterte weiter. Am Horizont löste sich
die Linie der Bergkette vollkommen auf, zerrann in einem milchigen,
grellen Band aus ungefiltertem Licht. Sie saßen da, während der Schatten
der Hütte sich langsam über den Boden schob, ein dunkler, scharfer Keil,
der sie schließlich beide verschlang, ohne ihnen Kühlung zu bringen.

Sie sprachen nicht mehr. Das Schweigen war nun vollständig, lückenlos,
versiegelt durch die unbewegliche Präsenz des Gesteins.

\section{\texorpdfstring{\textbf{Kapitel 10: Die Negation des
Umkreises}}{Kapitel 10: Die Negation des Umkreises}}\label{kapitel-10-die-negation-des-umkreises}

Am einundzwanzigsten Tag war die Hitze spiegelnd.

Sie stand nicht über der Senke, sie lag darin, fest und homogen wie ein
gegossener Block aus Silizium. Luca trat an die Türkante. Das Blech des
Oldsmobiles reflektierte das Mittagslicht in stumpfen, weißlichen
Kaskaden. Die Motorhaube wölbte sich leicht, von der Ausdehnung der
inneren Bauteile leicht verzogen.

Phillips saß auf der untersten Stufe. Seine Ellbogen ruhten auf den
Oberschenkeln, die Hände ineinandergefügt zu einer starren,
feingliedrigen Symmetrie. Seine Haltung war die eines Mannes, der eine
Berechnungsreihe abwartet, deren Laufzeit er nicht mehr beeinflussen
kann.

„Der Tank ist halb voll``, sagte Luca.

Das Wort \emph{Tank} verpuffte gegen das Wellblech.

Phillips blickte nicht auf. Ein trockenes, rhythmisches Geräusch kam aus
seinem Brustkorb, das Ausströmen von Luft durch verengte Atemwege,
präzise wie das Takten eines Kompressors. „Die Straße nach Westen
schneidet drei Sektoren. Es gibt dort keine Erfassungspunkte mehr. InSIM
hat die Masten vor vier Wochen abgebaut.``

Luca stieg in den Wagen. Die Lederpolster schickten eine stechende Welle
aus chemischer Hitze und synthetischem Harz durch das Chassis. Phillips
folgte ihm nach einer halben Minute. Er setzte sich auf den
Beifahrersitz, ohne die Tür zuzuschlagen; der Verschluss rastete erst
beim zweiten Versuchsstoß ein, ein trockenes Klacken von Eisen auf
Eisen.

Sie fuhren an. Die Räder griffen in den Bimsstein, warfen eine flache,
graue Salve gegen die Radkästen und fanden die verdichtete Spur des
Schotterwegs.

Die erste Ansiedlung war eine Ansammlung von geometrischen Körpern aus
Zement und Welleisen. Eine Tankstelle mit zwei verrosteten Zapfsäulen,
deren mechanische Zählwerke auf Null blockiert waren. Daneben ein
flacher Bau mit dem verblichenen Neonschriftzug eines Diners.

Luca ließ den Motor laufen. Die Erschütterung des V8-Blocks ging durch
das Lenkrad direkt in seine Handflächen über.

Er betrat den Raum. Die Luft roch nach Frittierfett, verbranntem
Filterkaffee und dem Ozon eines alten Fliegenfängers, der an der Decke
summte. Hinter der Resopal-Theke stand ein Mann mit einem blauen
Baumwollhemd, dessen Stoff an den Achseln dunkel verfärbt war. Seine
Pupillen bewegten sich nicht, als Luca an den Tresen trat.

Luca legte das Foto auf die Formica-Platte. Es war ein Computerausdruck
aus dem Archiv in Neapel, dreihundert DPI, das Gesicht leicht
rasteriert, die Kanten des Kiefers pixelig.

Der Mann sah nicht hin. Er blickte durch das Papier hindurch auf den
Metallsockel der Kaffeemaschine.

„Sie trägt eine dunkle Jacke``, sagte Luca. „Aus festem Baumwollstoff.
Das Haar zurückgebunden.``

„Wir haben hier keinen Durchgang``, sagte der Mann. Seine Stimme war
holzig, gereinigt von jeder lokalen Färbung. „Die Ausfallstraße ist seit
dem Zwölften gesperrt. Bauarbeiten am Umspannwerk.``

„Sie ist zu Fuß unterwegs.``

„Niemand geht zu Fuß über die Schwelle``, sagte der Mann. Er nahm einen
feuchten Lappen und wischte eine unsichtbare Ellipsenform in das
Resopal. Der nasse Film verdampfte augenblicklich unter den
Leuchtstoffröhren. „Es gibt hier keine Daten zu Fußgängern.``

Luca ließ das Bild liegen. Als er die Glasflügeltür aufdrückte, schwang
sie stumm hinter ihm zu, ein Vakuumverschluss, der die Hitze des
Außenraums drinnen hielt.

Sie erreichten die Kreuzung drei Meilen weiter südlich.

Dort stand kein Gebäude mehr. Nur ein Betonfundament, aus dem vier
verrostete Gewindestangen in den Himmel ragten. Der Anblick ähnelte den
Grundmauern der Villa dei Misteri, wo die Archäologen im neunzehnten
Jahrhundert den Gips in die Hohlräume gegossen hatten, um die Form des
verbrannten Holzes zu retten.

Luca hielt an. Er schaltete den Motor nicht aus. Der Auspuff spuckte
eine unregelmäßige, bläuliche Abgasfahne über den heißen Schotter.

Phillips blickte starr durch die Windschutzscheibe. Das Glas war von
mikroskopisch kleinen Einschlägen übersät, ein Fleckenmuster aus
gebrochenem Licht.

„Es gibt keine Martina in den Außenbezirken``, sagte der alte Mann.
Seine Stimme klang wie ein vorgelesenes Protokoll. „Du suchst nach einer
Ortsangabe in einer Matrix, die keine Raumkoordinaten mehr unterstützt.
Sie hat den Vektor gelöscht, Luca. Nicht den Körper.``

„Ein Körper hinterlässt Spuren``, sagte Luca. Er umklammerte das
Kunststoff-Lenkrad. „Er verdrängt Masse. Er verbraucht Wasser.``

„In der Iteration von Pompeji gab es vierhundert Leichname, von denen
man nur die Sandalen gefunden hat``, antwortete Phillips. Er bewegte den
Kopf nicht. „Kein Knochenmehl, keine Zähne. Nur die eiserne Schnalle im
Vulkanit. Weil das Signal abgerissen ist, bevor die Asche abkühlte. Das
ARS tut nichts anderes. Es vervollständigt die Lücke mit der
wahrscheinlichsten Hypothese. Und die wahrscheinlichste Hypothese für
diese Senke ist: Nichts.``

Luca legte den Rückwärtsgang ein. Das Getriebe heulte kurz auf. Der
Wagen drehte auf dem engen Schotterstreifen, die Reifen fraßen sich tief
in den Bims, bis sie wieder die Gegenfahrbahn erreichten.

Als sie zur Hütte zurückkehrten, war der Schatten der Veranda auf wenige
Zentimeter geschrumpft.

Luca ging direkt in die hintere Kammer. Das Notizbuch lag noch immer auf
der Fichtenplatte, exakt parallel zur Tischkante. Die Ecken des
Wachstuchs schienen durch die Wärmeeinwirkung leicht verzogen.

Er schlug die Seite auf.

Die Zeile am oberen Rand stand unberührt im Raum:

\emph{In einer anderen Iteration habe ich die Parameter nicht
bestätigt.}

Er fuhr mit dem Zeigefinger über die Tinte. Sie war vollkommen trocken,
glatt, eingebrannt in die Fasern wie die Rußschichten auf den Fresken
der Regio V.

Er trat hinaus auf die Veranda. Phillips saß wieder am selben Fleck, die
Fersen fest in den heißen Staub gestemmt.

„Sie wusste von den Protokollen``, sagte Luca.

Phillips schloss für zwei Sekunden die Augen. Das Leder seiner Lider war
dünn und trüb. „Sie hat die Protokolle geschrieben. Die
Dialoggrammatiken für die Sprachausgabe von InSIM stammten von ihr. Sie
wusste, dass das System nur Antworten verarbeitet, die eine logische
Bestätigung enthalten. Eine Verneinung ist eine Eingabe. Eine
Verweigerung ist das Rauschen zwischen zwei Abfragen.``

„Sie hat sich gelöscht.``

„Sie hat die Schleife unterbrochen``, korrigierte Phillips leise. Er
zeigte mit dem abgemagerten Zeigefinger auf die weiße Fläche der Senke.
„Wenn Teilhard de Chardin vom Punkt Omega sprach, meinte er die finale
Verdichtung des Bewusstseins. Er lag falsch. Es gibt keine Verdichtung.
Es gibt nur die unendliche Ausdünnung der Wahrscheinlichkeiten. Wir
sitzen hier im kleinsten gemeinsamen Nenner der Datenmenge.``

„Wir sind hier``, sagte Luca. „Ich spüre die Hitze. Ich rieche das
Benzin.``

„Das sind Sensordaten``, sagte der alte Mann. „Synthetische
Rückmeldungen einer fehlerhaften Kalibrierung. Du glaubst, du fährst
eine Straße entlang, weil dein Gehirn die Abfolge von Bildern verlangt.
Aber der Tank ist nicht leer geworden, Luca. Der Zeiger steht seit drei
Tagen auf derselben Markierung.``

Luca blickte auf seine Handflächen. Die Haut war trocken, von einer
feinen, grauen Staubschicht überzogen, die sich nicht abreiben ließ.

Er setzte sich auf die Stufe neben den Jesuiten.

Das Licht bröckelte langsam an den Rändern der Bergkette ab. Kein
gewöhnlicher Sonnenuntergang mit roten Farbsäumen, sondern ein
stufenweises Herunterregeln der Helligkeit, als würde die Ausleuchtung
der Senke in diskreten Quantensprüngen reduziert.

„Was bleibt übrig?{\kern0pt}``, fragte Luca.

„Das Material``, sagte Phillips. Seine Stimme war nun kaum mehr als ein
Flüstern, das im Knirschen des abkühlenden Wellblechs unterging. „Die
ungemessene Masse. Die Steine, die keine Signale mehr senden.``

Sie blieben sitzen. Der Schatten schob sich über ihre Knie, kühllos und
schwer wie flüssiges Blei, und löschte die Konturen ihrer Körper im
Dämmerlicht aus.

\section{\texorpdfstring{\hfill\break
}{ }}\label{section-4}

\section{\texorpdfstring{\textbf{Kapitel 11: Die Interferenzen des
Beobachters}}{Kapitel 11: Die Interferenzen des Beobachters}}\label{kapitel-11-die-interferenzen-des-beobachters}

Am dreiundzwanzigsten Tag war die Nacht ein fester Körper.

Luca wachte nicht durch ein Geräusch auf, sondern durch das Fehlen der
thermischen Spannung im Gebälk. Die Hütte war ausgekühlt, ein Raum ohne
Resonanz. Die Wachskerze auf der Konsolentischplatte war bis auf den
Docht heruntergebrannt; der geschmolzene Talg war starr und weiß wie die
Kalkschichten in den Kanälen von Herculaneum.

Er trat auf die Veranda.

Der Mond stand im Zenit, ein flacher, silbriger Schnittpunkt im stumpfen
Blau des Himmels. Die Bimssteinsenke reflektierte das Licht in einer
harten, schattenlosen Intensität. Jeder Vulkanitblock besaß die Schärfe
einer mathematischen Projektion, als liege die Senke nicht im Raum,
sondern in einem gerenderten Testfeld von InSIM.

Er sah die Gestalt am unteren Ende des Schotterstreifens.

Sie bewegte sich ohne das typische Schwanken eines Körpers, der Masse
über unebenen Boden verlagert. Es war ein gleichmäßiges Gleiten über das
Sediment.

Der Mann blieb drei Meter vor den Holzstufen stehen. Er trug eine
dunkle, taillierte Jacke aus schwerer Baumwolle -- denselben Schnitt,
den Phillips auf den Aufnahmen des Konsortiums von 1994 getragen hatte,
als die ersten Algorithmen für die Archäometrie von Sektor V programmed
worden waren.

Sein Gesicht war unbeschädigt von der UV-Strahlung der Senke. Keine
tiefen Kerben an den Schläfen, keine Hornhautmembran über den Pupillen.
Es war die Physiognomie des Jesuiten aus dem Zeitungsausschnitt im
Archiv von Neapel: das schmale, akademische Kinn, die steile Stirnfalte
eines Mannes, der stundenlang über Matrizen brütete.

„Die Wahrscheinlichkeitsdichte in diesem Sektor ist extrem niedrig``,
sagte der Mann. Seine Stimme hatte keinen Nachhall, als würde der Schall
direkt von der Oberfläche der Bimssteinblöcke absorbiert. „Ihr solltet
den Brenner ausschalten. Das Abgas verändert die Isotopenwerte im
Sediment.``

Luca richtete sich nicht auf. Seine Finger umfassten die raue, vom Harz
klebrige Kante des Geländers. „Wir haben dich in den Außenbezirken
gesucht.``

„Dort gibt es keine Überlagerungen mehr``, sagte der Jüngere. Er blickte
nicht zu Luca, sondern über das Wellblechdach hinweg auf die Linie der
Hügelkette. „InSIM hat die Abtastfrequenz gesenkt. Was ihr Außenbezirke
nennt, ist nur noch der Pufferbereich zwischen zwei Iterationen. Wenn
eine Variable ihre Rückmeldung verweigert, zieht das ARS den
Parameterbereich zusammen. Martina ist nicht gegangen. Sie hat nur die
Kohärenz mit dieser Speichereinheit verloren.``

„Das ist ein theoretisches Konstrukt``, sagte Luca.

„Es ist die Architektur, in der du sitzt``, antwortete der Mann. Er hob
die Hand, eine kurze, fast beiläufige Bewegung. Die Haut an seinen
Knöcheln war glatt, ohne Flecken. „Phillips hat in den siebziger Jahren
in Rom über die Logik der Verschränkung publiziert. Er wusste, dass man
ein Bewusstsein nicht simulieren kann, ohne ihm den Austritt aus dem
System zu erlauben. Wenn der Beobachter die Messung nicht mehr
bestätigt, kollabiert der Zustand nicht. Er verweilt in der
Überlagerung.``

„Wer bist du?{\kern0pt}``, fragte Luca.

Der Mann wandte das Gesicht langsam der Veranda zu. Seine Augen fingen
das Mondlicht wie zwei polierte Linsen. „Der Zustand vor der
Entscheidung. Der Teil des Algorithmus, der in den Archiven von Pompeji
geblieben ist, als das Projekt die ersten Freigaben erhielt. Wir haben
das Modell nicht gebaut, um die Toten wiederzubeleben, Luca. Wir haben
es gebaut, um die Notwendigkeit einer einzigen Realität aufzuheben.``

„Phillips liegt im Zimmer``, sagte Luca. „Er stöhnt im Schlaf. Seine
Venen treten an den Beinen hervor.``

„Er stirbt in der Hauptschleife``, sagte der jüngere Phillips ohne
Härte. „Das ist die Bedingung für diese Messreihe. Ein Beobachter, der
im Zerfall begriffen ist, erzeugt die schärfsten Signale. Aber Martina
benötigt das Signal nicht mehr. Sie befindet sich in einer Iteration,
deren Parameter nicht abgeglichen werden müssen.``

Er machte keine Kehrtwendung. Er trat lediglich einen Schritt zurück,
und mit diesem Schritt schien der Kontrast seiner Gestalt zu schwinden,
als verliere die Auflösung seiner Konturen an Tiefe. Die dunkle
Baumwolle seiner Jacke floss in das Grau der Bimssteine über.

„Wenn du weiter nach ihr suchst``, sagte die Stimme, die nun aus der
Richtung der Felszunge zu kommen schien, „verstärkst du nur das
Rauschen. Das ARS verarbeitet keine Trauer. Es verarbeitet nur
Koordinaten.``

Dann war das Sediment leer. Kein Fußabdruck, kein verlagertes
Kieselgeröll. Nur das flache, monotone Flimmern des Mondlichts auf dem
Basalt.

Am vierundzwanzigsten Tag stand die Sonne wie ein weißglühender Stempel
über dem Wellblech.

Luca saß auf den Holzstufen. Der Staub auf seinen Stiefeln war durch die
Luftfeuchtigkeit der Nacht zu einer dünnen, graubeigen Kruste
getrocknet.

Phillips trat aus der Tür. Seine Bewegungen hatten die Brüchigkeit von
morschem Holz, das unter minimaler Belastung nachgibt. Er setzte sich
mit einem langen, pfeifenden Ausatmen neben Luca auf die oberste Stufe.
Das Pflaster an seinem Hals war an den Rändern braun verfärbt und löste
sich ab.

„Der Gradient fällt``, sagte Phillips. Er sah auf seine eigenen Hände,
die kraftlos auf den Kniescheiben ruhten. „Der Prozessor regelt die
Rechenleistung herunter. Hast du die Frequenz gehört?{\kern0pt}``

„Er war heute Nacht hier``, sagte Luca.

Phillips bewegte den Kopf nicht. „Die Rekonstruktion greift auf
Altbestände zurück. Wenn der primäre Speicher beschädigt ist, holt das
System frühere Zustände aus dem Zwischenspeicher. Das ist kein Geist,
Luca. Das ist ein Lese-Schreib-Fehler in den Registern.``

„Er trug deine Jacke aus Neapel.``

„Er trug die Textur, die das ARS für mein dreißigjähriges Ich hinterlegt
hat``, sagte der alte Mann. Seine Stimme war trocken, eine
Konsonantensequenz ohne Stimmbandvibration. „Das System versucht, die
Lücke zu schließen, die Martina hinterlassen hat. Es bietet dir ein
Ersatzsignal an, damit du den Beobachterstatus nicht aufgibst. Solange
du antwortest, läuft die Simulation weiter.``

Luca griff in den Staub zwischen seinen Stiefeln und hob einen Brocken
Vulkanit auf. Der Stein war leicht, durchsetzt mit winzigen Hohlräumen,
gefangenem Gas aus einem Ausbruch, der zweitausend Jahre zurücklag.

„Und wenn wir aufhören zu antworten?{\kern0pt}``, fragte Luca.

Phillips sah auf den Stein in Lucas Hand. Ein minimales, zitterndes
Licht spiegelte sich in seinen trüben Pupillen.

„Dann gibt es keine Iteration mehr``, sagte der Jesuit. „Kein Pompeji.
Keine Senke. Nur die ungemessene Masse. Die reine, zeitlose Gegenwart
des Gesteins.``

Sie saßen nebeneinander auf der Stufe. Die Hitze stieg aus dem Bimsstein
auf wie das Flimmern über einer frischen Ascheschicht. Der Schatten der
Veranda verkürzte sich Zentimeter um Zentimeter, bis er als schmaler,
dunkler Strich unter dem Holz verschwand und sie vollkommen dem
ungefilterten Licht überließ.

\section{\texorpdfstring{\hfill\break
}{ }}\label{section-5}

\section{\texorpdfstring{\textbf{Kapitel 12: Das Ende des
Beobachters}}{Kapitel 12: Das Ende des Beobachters}}\label{kapitel-12-das-ende-des-beobachters}

Am fünfundzwanzigsten Tag schrumpfte das Licht auf eine thermische
Frequenz.

Luca erwachte mit der Wange auf der rauen Kiefernholzplatte des Tisches.
Unter seiner Schläfe lag das Oktavheft, dessen Papier durch die extreme
Trockenheit der Wüstenluft spürbar vergilbt war. Die Ränder der Seiten
rollten sich auf wie getrockneter Pflanzensaft.

Phillips saß auf der Veranda. Er war bewegungslos, gebückt über seine
eigenen Oberschenkel, als beuge er sich über ein kaum wahrnehmbares
Schwingungsmuster im Sediment. Das Pflaster an seinem Hals hatte sich an
der linken Ecke gelöst; darunter lag die blasse, narbige Spur einer
alten Einstichstelle, grau wie die Kalziumkrusten an den Brunnenwänden
von Pompeji.

Luca nahm den Füllfederhalter auf. Die Tinte floss zäh, unterbrochen von
Mikro-Aussetzern im Fluss. Er notierte nicht mehr die Abfolge der
Stunden, sondern die physikalische Ablagerung von Zeit: die Ausdehnung
des Bimssteins unter der Mittagshitze, das tonlose Erkalten der
Metallklammern im Dachstuhl, die Überlagerungszustände, die das ARS als
Rauschen durch die Außenbezirke schickte.

Er schrieb die Syntax auf, die der Jesuit in Neapel konzipiert hatte.
Die Verschränkung von Materie und Gedächtnis. Das Modell Teilhard de
Chardins, übersetzt in die binäre Stille eines Quantenprozessors: der
Omega-Punkt nicht als Erlösung, sondern als maximale Verdichtung aller
nicht gewählten Iterationen. Martina war nicht fortgegangen. Sie hatte
lediglich den Beobachtungsraum verlassen, um im ungemessenen Feld zu
verbleiben.

In der Nacht trat Phillips in den Raum. Er schleifte die Fersen über die
Dielen -- ein trockenes, schabendes Geräusch wie Treibsand auf Blech. Er
legte sich auf das Feldbett, ohne die Schuhe auszuziehen. Sein Atem war
flach, synchronisiert mit dem periodischen Summen des Solargenerators
hinter der Hütte.

Am siebenundzwanzigsten Tag stellte InSIM die Signalverarbeitung ein.

Der Flimmerstreifen am Horizont, der über Wochen die Grenze zwischen der
Senke und den Zementabbaugebieten markiert hatte, war verschwunden. Die
Landschaft hatte ihre Tiefe verloren; die Felsformationen wirkten flach,
eingestanzte Silhouetten gegen einen Himmel aus ausgebleichtem Zink.

Luca trat zu Phillips auf die Stufen.

„Die Schnittstelle bricht ab``, sagte Luca.

Phillips blickte nicht auf. Seine Finger nestelten an einem Fadenrest
seiner Jacke. „Es gibt keine Schnittstelle mehr. Das Konsortium hat den
Speicherbereich freigegeben. Was du siehst, ist der Ruhezustand der
Matrix vor der ersten Datenaufnahme.``

„Das Oktavheft ist voll``, sagte Luca.

„Das Heft ist nur eine zusätzliche Speichereinheit``, antwortete der
Jesuit. Seine Stimme war kaum mehr als ein zischender Luftstrom, der an
den ausgezehrten Lippen brach. „Du hast die Daten interpretiert. Das war
dein Fehler. Das ARS zeichnet keine Bedeutungen auf. Es registriert
lediglich die Dichte von Zuständen. Wenn du versuchst, eine Bedeutung zu
finden, erzeugst du ein Rauschen, das das Signal zerstört.``

„Ich bin morgen weg``, sagte Luca.

Phillips schwieg. Ein Schatten flog über die Bimssteinsenke -- kein
Vogel, sondern die schemenhafte Verschiebung einer Pixellinie im oberen
Aufnahmefeld.

„Der Wagen steht am Rand der Schotterpiste``, sagte Luca. „Die
Batterieladung reicht bis an die Küste.``

„Es gibt keine Küste mehr``, sagte Phillips leise. Er schloss die Augen.
Die Lider waren so dünn, dass die feinen, bläulichen Verästelungen der
Äderchen darunter sichtbar blieben wie die Mikro-Fissuren in der
Ascheschicht von Sektor VI. „Es gibt nur das Sediment und den
Beobachter, der den Zerfall mitschreibt.``

Am achtundzwanzigsten Tag war der Raum leer von Schatten.

Die Sonne stand senkrecht über der Senke, ein weißer, strahlungsloser
Punkt im Zentrum der Wahrnehmung.

Luca packte die Reisetasche. Das Gewicht des Oktavhefts lag schwer am
Boden der Quertasche, ein rechteckiger Block aus Zellulose und trockener
Tinte. Er trat hinaus auf die Veranda.

Phillips saß an der Wand angelehnt. Seine Augen waren geöffnet, aber der
Blick fixierte keinen Punkt der Hügelkette mehr; die Pupillen waren
starr, reflektierten das monotone Grau der Bimssteinblöcke wie zwei matt
geschliffene Basaltlinsen. Seine Brust bewegte sich nicht mehr.

Luca setzte sich nicht. Er blieb am Fuß der Treppe stehen. Der Staub um
seine Stiefel war reglos, frei von thermischem Auftrieb.

„Die Messung ist abgeschlossen``, sagte Luca in die Stille.

Keine Antwort. Die Textur der Jacke des Alten schien mit dem rauen Holz
der Hüttenwand zu verschmelzen, eine graduelle Nivellierung der
Oberflächenspannungen.

Luca ging zum Fahrzeug. Der Motor startete ohne Mechanikgeräusch, ein
tiefer, niederfrequenter Impuls, der im Sediment abstarb. Er legte den
Rückwärtsgang ein und wandte den Wagen auf dem schmalen
Schotterstreifen.

Im Rückspiegel schrumpfte die Hütte zu einem flachen Quadrat aus
Wellblech und Holz. Die Gestalt auf der Veranda verlor ihre Konturen,
floss in die grauen Farbtöne der Blöcke über, bis nur noch der Hang der
Senke sichtbar blieb.

Er fuhr durch die Pufferzone. Die Straße bildete sich erst wenige Meter
vor den Scheinwerfern aus -- harter, abgetragener Asphalt mit den gelben
Trennlinien des Highway 190. Nach vierzig Meilen erschienen die ersten
Hochspannungsmasten, abstrakte Stahlgitterstrukturen gegen das blasse
Reorganisationslicht des Nachmittags.

In Los Angeles war der Flughafen ein Raum aus Milchglas und gefilterter
Kühlluft.

Luca saß am Gate 42. Durch die Panoramascheibe sah er die Rumpfflächen
der Maschinen, glänzend, ohne Nieten, saubere Geometrien in einem
Kontinuum aus Kerosingeruch und Lautsprecherdurchsagen. Die Passagiere,
die an ihm vorbeigingen, trugen Gesichter ohne Tiefenschärfe, flüchtige
Rekonstruktionen in einer dicht besetzten Schleife.

Er öffnete die Quertasche und zog das Oktavheft heraus.

Er schlug die erste Seite auf. Das Papier war leer. Keine
Tintenstreifen, keine Eintragungen, keine Diagramme über die
Verschränkung von Impuls und Ort. Nur die gleichmäßige, jungfräuliche
Struktur des Papiers.

Er blätterte zurück. Seite für Seite. Das gesamte Heft war weiß. Die
Zeichen hatten sich aufgelöst, als hätte das ARS die nicht mehr
benötigten Parameter nach dem Austritt aus der Senke überschrieben.

Die Stimme aus den Deckenlautsprechern rief die Gruppe B auf.

Luca stand auf. Er schob das leere Heft zurück in die Tasche, schloss
den Reißverschluss und trat in das Fingerdock. Der Passagierraum des
Flugzeugs war kühl und roch nach Ozon. Er nahm den Fensterplatz ein,
schnallte den Gurt fest und sah zu, wie die Maschine vom Schlepper
zurückgesetzt wurde.

Die Startbahn schob sich unter den Flügeln durch. Beim Abheben kippte
die Landschaft leicht zur Seite: das Häusermeer von Los Angeles, das
Raster der Freeways, die dünne Linie der Küste -- alles zog sich
zusammen zu einer flachen, hochauflösenden Textur.

Er schloss die Augen.

In der Senke veränderte sich die Helligkeit nicht. Der Bimsstein
absorbierte das verbliebene Licht. Die Blöcke lagen im Staub, ohne
Schattenwurf, ohne Temperaturgradienten, ungemessen und zeitlos.

\section{\texorpdfstring{\hfill\break
}{ }}\label{section-6}

\section{Kapitel 13: Die Rückkehr}\label{kapitel-13-die-ruxfcckkehr}

Die Maschine setzte um sechzehn Uhr dreiundzwanzig auf der Landebahn in
Boston auf. Das Fahrwerk stieß mit einem kurzen, trockenen Schlag gegen
die Betonplatten, dann rollte der Rumpf aus, und die Triebwerke schoben
die Verzögerung in die Sitzlehnen. Luca saß am Fenster. Er hatte während
des gesamten Fluges nicht geschlafen. Vor ihm auf dem Klapptisch stand
eine unberührte Plastiktasse mit Wasser, deren Oberfläche noch die
feinen Rillen des Einschenkens trug.

Draußen zog die Silhouette der Stadt vorbei. Backsteinfronten,
Lagerhallen, das stumpfe Grau des Hafens. Boston war nicht Los Angeles.
Boston war nicht die Wüste. Boston war der Ort, an dem er hergekommen
war, und der Ort, an den er zurückkehrte, aber nichts davon fühlte sich
wie Heimkehr an.

Im Flugzeug hatte er das Oktavheft geöffnet. Die Seiten waren leer. Er
hatte sie zweimal durchgeblättert, Seite für Seite, in der Erwartung,
dass irgendwo zwischen den Blättern ein Rest Tinte geblieben war, ein
Diagramm, eine Randnotiz. Nichts. Nur das reine, jungfräuliche Papier.
Das ARS hatte die Parameter überschrieben, als er die Senke verließ.
Nicht gelöscht. Ersetzt. Das Heft war kein Heft mehr. Es war der
Grundzustand eines Heftes.

Er schob es in die Quertasche, ohne es zu schließen.

Der Weg vom Logan Airport in die Innenstadt dauerte länger als sonst.
Die Silver Line stand im Stau vor der Ted-Williams-Tunnel-Einfahrt. Die
Fahrgäste im Bus starrten auf ihre Telefone, deren Displays in
gleichmäßigem Rhythmus aufleuchteten. Niemand sprach. Luca stand, die
Hand an der Haltestange, und sah durch die getönte Scheibe auf die
Straße. Die Ampeln schalteten im Sekundentakt. Die Menschen gingen. Die
Autos hielten. Die Stadt funktionierte.

Er stieg an der South Station aus und ging zu Fuß weiter. Die Luft war
kühl, feucht, nach Salz und altem Regen. Anders als die trockene,
alkalische Luft der Senke. Die Haut in seinem Gesicht, die von der Wüste
rissig geworden war, zog sich in der Feuchtigkeit zusammen.

Die Wohnung im vierten Stock roch noch nach dem, was sie immer gerochen
hatte: Kaffee, Staub, altes Holz. Seine Mutter war nicht da. Auf dem
Küchentisch lag ein Zettel, mit ihrer Handschrift: Bin bei Tante Ruth.
Komm später. Es gibt Suppe. Luca las den Zettel zweimal, dann legte er
ihn zurück auf die Granitimitat-Arbeitsplatte.

Er duschte. Er wechselte die Kleidung. Er setzte sich auf das Sofa und
starrte an die Wand. Dann stand er wieder auf, zog die Jacke an und ging
hinaus.

Der Seminarraum lag im dritten Stock des Kommunikationsgebäudes. Luca
hatte nicht gewusst, ob der Dozent an diesem Dienstag unterrichtete. Er
hatte es auch nicht geprüft. Er war einfach gegangen, in der Annahme,
dass der Raum da war, das Licht brannte, die Jalousien ihre Streifen
warfen.

Er hatte recht.

Der Raum war halb gefüllt. Zehn, zwölf Studierende in den Reihen, die
Köpfe über Notizbücher gebeugt. An der Stirnwand flimmerte die
Projektion. Auf dem Bildschirm lief Mr. Nobody, dieselbe Szene wie vor
Wochen, auf dieselbe Langsamkeit heruntergedimmt. Der neunjährige Nemo
auf dem Bahnsteig. Die Mutter mit dem Koffer. Der Vater neben dem
Jungen. Die beiden Blicke. Der Zug. Der Schnürsenkel.

Der Dozent stand seitlich des Strahlers. Die Brille spiegelte das
Flimmern. Er sagte nichts. Er ließ das Bild arbeiten.

Luca blieb an der Tür stehen. Er sah auf den Bildschirm, dann auf die
Gesichter der Studierenden, dann wieder auf den Bildschirm. Nichts hatte
sich verändert. Alles hatte sich verändert.

Er ging in die dritte Reihe und setzte sich auf den Platz, auf dem er
vor Wochen gesessen hatte. Die Unterarme auf die Kunstharzplatte. Der
Pullover war derselbe graue. Die Ellbogen lagen an derselben Stelle.

Der Dozent bemerkte ihn nicht sofort. Als er es tat, unterbrach er
nicht. Er ließ die Szene weiterlaufen. Die Kamera hielt beide Eltern im
Kader, während die Unschärfe wechselte. Der Junge blickte von der Mutter
zum Vater und zurück. Der Schnürsenkel riss. Das Kind stolperte. Der Zug
fuhr aus dem Bild.

Dann sagte der Dozent den Satz, den Luca schon gehört hatte: „Wir
sprechen hier nicht von einer einfachen Dramaturgie der Wahl."

Luca hörte die Worte, aber sie klangen anders. Sie klangen flacher. Sie
klangen nach etwas, das nicht mehr reichte.

Der Dozent fuhr fort. Er sprach von der Quanten-Schnittstelle, von den
probabilistischen Pfaden, von der Rechenarchitektur der
Viele-Welten-Interpretation. Er sprach von Teilhard und vom Omega-Punkt
als der ultimativen Verdichtung des Bewusstseins. Er stellte die Frage,
was passiere, wenn die Rechenleistung eines Systems diese Singularität
auflöse und jede Möglichkeit, jedes Stolpern am Bahnsteig, simultan real
halte. Wo bleibe die Schuld, wo bleibe der Wille, wenn jede Kontingenz
eine neue Welt erzeuge.

Dieselben Fragen wie vor Wochen. Dieselben Antworten. Dieselbe Sprache.

Luca hob die Hand.

Der Dozent sah ihn an. Ein kurzes Zögern, dann ein Nicken.

„Ja?"

Luca stand nicht auf. Er blieb sitzen, die Hände auf der Tischplatte,
den Blick auf den Bildschirm gerichtet, wo die Szene in einer
Wiederholungsschleife weiterlief.

„Sie haben letztes Mal gesagt, dass Sprachmodelle und stochastische
Prozesse keine Urheber liefern", sagte Luca. „Nur Kohärenzen."

Der Dozent schwieg. Die Brille spiegelte das Flimmern.

„Ich habe darüber nachgedacht", fuhr Luca fort. „In der Wüste. In einer
Senke. Bei einem alten Mann, der die Sprache abgelegt hat. Und ich
glaube, Sie haben recht, aber nicht auf die Art, wie Sie es gemeint
haben."

„Wie meinen Sie es?", fragte der Dozent.

Luca sah auf den Bildschirm. Der Junge am Bahnsteig. Die Mutter. Der
Vater. Der Zug, der aus dem Bild fährt. Der Schnürsenkel, der reißt.

„Es gibt zwei Arten von Kohärenz", sagte Luca. „Die eine ist die von
DeLillo. Die andere ist die von Van Dormael."

Der Dozent trat einen halben Schritt zur Seite. Sein Gesicht war nun
voll im Licht. Die Falten um seine Augen traten schärfer hervor. Er
wartete.

„DeLillo zeigt die Kohärenz, die zerbricht", sagte Luca. „Die Stille.
Die Wüste. Die Steine. Alles läuft auf einen Grundzustand zu, in dem
keine Unterscheidung mehr möglich ist. Kein Urheber. Kein Wille. Keine
Schuld. Nur das ungemessene Material. Das ist der Omegapunkt als
Wärmetod. Als Entropie. Als Ausdünnung."

Er sah auf die Studierenden. Einige hatten die Stifte abgelegt. Andere
sahen noch auf ihre Notizbücher.

„Van Dormael zeigt die andere Art", sagte Luca. „Die Kohärenz, die sich
vervielfacht. Die Verzweigungen, die nicht ineinander verschwinden. Die
Zeit als Feld, nicht als Strom. Jede Entscheidung eine Welt, jede Welt
ein Weg, jeder Weg der richtige Weg. Das ist der Omegapunkt als Fülle.
Als Komplexität der Möglichkeiten. Als Differenz, die nicht nivelliert
wird."

Er hielt inne. Der Bildschirm hinter ihm zeigte weiter die Schleife. Der
Junge am Bahnsteig. Der Schnürsenkel. Der Zug.

„Beide Perspektiven sind richtig", sagte Luca. „Aber sie sind nicht
vereinbar. DeLillo und Van Dormael lassen sich nicht zu einer Synthese
bringen. Sie sind zwei Perspektiven auf denselben Grenzwert. Der eine
sieht das Ende. Der andere sieht die Fülle."

Der Dozent schwieg. Die Studierenden schwiegen. Nur die Projektion lief
weiter.

„Und Teilhard?", fragte der Dozent schließlich.

„Teilhard", sagte Luca, „hat eine dieser Perspektiven für die ganze
gehalten. Die teleologische. Die mit dem Ziel. Die mit der Richtung. Er
war Augustiner. Er war Thomist. Er konnte nicht anders, als die Zeit auf
ein Ende hin zu denken. Das war seine Größe. Und sein Irrtum."

Ein kurzes, kaum hörbares Atmen ging durch die Reihen. Der Dozent sah
Luca an, ohne den Kopf zu bewegen. Die Brille fing einen neuen
Lichtstreifen auf.

„Sie haben in der Wüste etwas gesehen", sagte der Dozent.

„Ja", sagte Luca.

„Was?"

Luca sah auf seine Hände. Die Haut war noch rissig von der Trockenheit.
Die Knöchel waren weiß. Er dachte an Phillips auf der Veranda. An den
jüngeren Phillips im Trümmerfeld. An die beiden Gestalten, die
gleichzeitig da waren und nicht vereinbar. An die Steine, die nicht
zeugten. An die Möglichkeiten, die nicht verschwanden.

„Ich habe gesehen, dass beide Perspektiven gleichzeitig gelten", sagte
er. „Nicht abwechselnd. Nicht nacheinander. Nicht als Kompromiss.
Sondern als Rand. Als Grenzwert. Als das, was übrig bleibt, wenn man
aufhört, sich zu entscheiden."

Der Dozent trat einen Schritt nach vorn. Er nahm die Brille ab und rieb
die Gläser mit einem Zipfel seines Hemdes. Ohne die Brille wirkten seine
Augen kleiner, verletzlicher. Er setzte sie wieder auf.

„Und was folgt daraus?", fragte er.

Luca zögerte. Nicht, weil er die Antwort nicht wusste. Sondern weil er
sie nicht zu laut sagen wollte.

„Dass es keinen Grund gibt, sich für eine der beiden zu entscheiden",
sagte er. „Dass der Versuch, sie zu vereinbaren, sie zerstört. Dass der
Versuch, eine für die ganze zu halten, zur Ideologie wird. Dass der
einzige ehrliche Umgang mit dem Omegapunkt darin besteht, beide
Perspektiven auszuhalten. Ohne Synthese. Ohne Versöhnung. Ohne Ausweg."

Er sah den Dozenten an. Die Studierenden. Die Projektion. Den Jungen am
Bahnsteig, der nicht wusste, ob er zur Mutter oder zum Vater gehen
sollte, und der am Ende weglief, weil er sich nicht entscheiden konnte.

„Das ist kein Nihilismus", sagte Luca. „Das ist die einzige Form, in der
man dem Leiden, der Zeit und der Endlichkeit gerecht werden kann, ohne
sie zu rechtfertigen. Kein Trost. Keine Erlösung von der Welt. Nur die
Vereinigung innerhalb des Gelebten."

Der Dozent sagte nichts. Er drehte sich zur Projektion um und ließ die
Szene von vorn beginnen. Der Junge am Bahnsteig. Die Mutter. Der Vater.
Der Schnürsenkel.

„Sie haben mehr gesagt, als Sie vielleicht wissen", sagte der Dozent,
ohne sich umzudrehen. „Sie haben die katholische Tradition nicht
verlassen. Sie haben sie nur von einer anderen Seite gelesen."

Luca stand auf. Er nahm seine Jacke. Er ging zur Tür. Er blieb auf der
Schwelle stehen.

„Das ist der Unterschied zwischen uns", sagte Luca. „Sie lesen die
Tradition. Ich habe sie gesehen."

Er ging.

Die Tür fiel hinter ihm ins Schloss. Der Flur war still. Die
Nachmittagssonne stand tief über den Dächern von Boston und warf lange,
schmale Streifen über den Steinboden. Er ging die Treppe hinunter, Stufe
für Stufe, im gleichen Rhythmus wie beim ersten Mal.

Draußen auf der Straße stand er einen Moment. Die Luft war kühl. Der
Himmel war grau. Die Autos fuhren. Die Menschen gingen. Die Ampeln
schalteten. Alles war da. Nichts war wichtig.

Er ging zu Fuß zurück. Durch die Straßen, die er kannte. An den Cafés
vorbei. An den Buchläden. An der Kirche, an der er als Kind
vorbeigegangen war. Die Kirchtür war offen. Ein Licht brannte drinnen.
Er blieb nicht stehen.

Er dachte an Phillips. An den alten Mann in der Wüste. An den jüngeren
Mann im Trümmerfeld. An die beiden Perspektiven. An den Rand, an dem man
aufhört, sich zu entscheiden.

Er dachte an Martina. An ihr Oktavheft. An den Satz: In einer anderen
Iteration habe ich die Parameter nicht bestätigt.

Er dachte an sich. An den Doktoranden, der in die Wüste gefahren war, um
einen alten Mann zu finden, und der einen Grenzwert gefunden hatte, der
nicht zu vereinbaren war.

Er dachte an den Abend. An die Suppe, die seine Mutter gekocht hatte. An
den Zettel auf dem Küchentisch. An Tante Ruth.

Und dann, zum ersten Mal seit Wochen, dachte er nichts mehr. Er ging
einfach. Die Straße hinunter. Zu sich nach Hause.

\section{Epilog: Die Steine}\label{epilog-die-steine}

\subsection{I. Die Wüste}\label{i.-die-wuxfcste}

Der Wind hatte sich gelegt. Die Senke lag unter einem Himmel, der weder
blau noch grau war, sondern die Farbe von ausgebleichtem Papier hatte.
Die Sonne stand tief, aber ihr Licht war matt, als käme es durch eine
dünne Schicht aus Staub, die sich seit Tagen nicht mehr gesenkt hatte.

Phillips saß auf der Veranda. Die Hände flach auf den Oberschenkeln. Die
Augen offen, aber auf keinen Punkt gerichtet. Das Pflaster an seinem
Hals hatte sich vollständig gelöst. Darunter lag die Haut, dünn wie
Pergament, grau wie die Kalkschichten an den Brunnenwänden von Pompeji.

Er atmete. Noch. Der Atem ging flach, kaum sichtbar unter dem Stoff des
ausgewaschenen Hemdes. Aber er ging.

Neben ihm auf der Stufe lag das Oktavheft. Nicht das seine. Das andere.
Das von Martina. Es lag aufgeschlagen da, die Seiten leer bis auf den
einen Satz im hinteren Register: In einer anderen Iteration habe ich die
Parameter nicht bestätigt.

Phillips sah nicht auf das Heft. Er sah auf die Steine.

Die Basaltblöcke lagen im Sediment wie eine Anordnung von Datenpunkten
auf einem dunklen Feld. Keine Schatten. Keine Konturen. Keine
Temperatur. Nur das reine, ungemessene Gewicht von Silikaten und
Erstarrung.

Das ARS hatte die Simulation nicht abgeschaltet. Es hatte sie
eingefroren. Es gab keine Updates mehr. Keine Parameter. Keine
Messungen. Es gab nur noch den Zustand, in dem keine Alternative mehr
berechnet werden musste. Der Grenzwert. Der Rand. Der Omega-Punkt.

Phillips bewegte die Lippen. Kein Ton kam heraus. Nur die Form von
Worten, die niemand mehr hören würde.

Dann wurde die Bewegung langsamer. Dann hörte sie auf. Die Hände blieben
auf den Oberschenkeln liegen. Die Augen blieben offen. Der Atem blieb
aus.

Die Steine veränderten ihre Farbe nicht. Sie absorbierten das Licht. Sie
waren da. Sie waren immer da gewesen. Sie sahen das Verschwinden. Aber
sie bezeugten es nicht.

Über dem Trümmerfeld flimmerte die Hitze. Dann flimmerte auch sie nicht
mehr. Dann war nichts mehr, was sich hätte bewegen können. Nur die
Stille, die keine Stille war, sondern der Grundzustand.

\subsection{II. Boston}\label{ii.-boston}

Das Kino lag in der Kneeland Street, ein alter Bau mit einem verbeulten
Vordach und einer Leuchtreklame, die seit Jahren nicht mehr repariert
worden war. Luca hatte den Ort zufällig gefunden, an einem Abend, an dem
er nicht nach Hause wollte. Es lief ein Film, den er kannte. Mr. Nobody.
In der Originalsprache. Ohne Untertitel.

Er saß in der letzten Reihe. Der Saal war halb leer. Ein paar Studenten,
ein älteres Paar, ein Mann mit einer Tüte Popcorn, der nicht aß, sondern
nur die Tüte hielt. Die Leinwand flimmerte. Der Vorspann lief.

Und dann war der Junge da. Auf dem Bahnsteig. Die Mutter mit dem Koffer.
Der Vater daneben. Die beiden Blicke. Der Zug. Der Schnürsenkel.

Luca sah die Szene nicht mehr wie vor Wochen. Er sah sie doppelt. Einmal
als das, was sie war: eine Entscheidung, die nicht getroffen werden
kann. Und einmal als das, was sie bedeutete: eine Verzweigung, die alle
Möglichkeiten enthält.

Der Junge rannte los. Der Schnürsenkel riss. Das Kind stolperte. Der Zug
fuhr aus dem Bild. Und dann, in der nächsten Einstellung, war der Junge
wieder da. Auf dem Bahnsteig. Die Mutter mit dem Koffer. Der Vater
daneben.

Die Schleife lief weiter. Nicht als Wiederholung. Als Möglichkeit.

Luca lehnte sich zurück. Die Sitze waren alt, das Leder rissig. Er
spürte die Kälte des Kinos in den Knochen. Aber er spürte auch etwas
anderes. Etwas, das er in der Wüste verloren geglaubt hatte.

Er sah nicht eine Welt. Er sah viele. Er sah den Jungen, der zur Mutter
ging. Den Jungen, der beim Vater blieb. Den Jungen, der weglief. Den
Jungen, der stehen blieb. Alle waren da. Alle waren real. Keine war
richtiger als die andere.

Das war Van Dormael. Das war die Fülle. Die Komplexität der
Möglichkeiten. Die Differenz, die nicht nivelliert wird.

Luca sah eine Stunde zu. Dann noch eine. Als der Film zu Ende war, blieb
er sitzen. Die Lichter gingen an. Die anderen Zuschauer standen auf und
gingen. Der Mann mit der Popcorn-Tüte blieb auch sitzen. Sie sahen sich
nicht an.

Luca nahm die Quertasche und ging hinaus.

Auf der Straße war es kühl. Die Lichter der Stadt hingen tief. Die
Menschen gingen an ihm vorbei, ohne ihn zu sehen. Die Autos fuhren. Die
Ampeln schalteten. Alles war da. Nichts war wichtig.

Er ging die Kneeland Street hinunter, dann rechts in die Washington
Street, dann über die Brücke. Die Donau war nicht da. Die
Freiheitsbrücke war nicht da. Das war Boston, nicht Budapest. Aber der
Fluss war da. Der Charles. Grau und still.

Er blieb auf der Brücke stehen. Unter ihm floss das Wasser. Über ihm
hing der Himmel. Neben ihm gingen die Menschen vorbei.

Und Luca dachte: Der Omegapunkt ist nicht das Ende. Er ist nicht die
Fülle. Er ist der Rand, an dem man aufhört, sich zu entscheiden.

Er dachte an Phillips. An den alten Mann in der Wüste, dessen Atem
aufgehört hatte. An den jüngeren Phillips im Trümmerfeld, dessen
Konturen sich aufgelöst hatten. An die beiden Perspektiven, die
gleichzeitig galten.

Er dachte an Martina. An ihr Oktavheft. An den Satz: In einer anderen
Iteration habe ich die Parameter nicht bestätigt.

Er dachte an sich. An den Doktoranden, der in die Wüste gefahren war, um
einen alten Mann zu finden, und der einen Grenzwert gefunden hatte, der
nicht zu vereinbaren war.

Er dachte an den Abend. An die Suppe. An den Zettel auf dem Küchentisch.
An seine Mutter.

Dann dachte er nichts mehr. Er ging weiter. Über die Brücke. In die
Stadt. Zu sich nach Hause.

\subsection{III. Die Straße}\label{iii.-die-strauxdfe}

Es war später Abend, als Luca vor dem Haus stand, in dem er aufgewachsen
war. Die Backsteinfassade. Die schmale Treppe. Das Licht im vierten
Stock, das noch brannte.

Seine Mutter war da. Sie hatte die Suppe warm gehalten. Sie fragte
nicht, wo er gewesen war. Sie fragte nicht, was er gesehen hatte. Sie
stellte ihm einen Teller hin und setzte sich ihm gegenüber.

Luca aß. Die Suppe war heiß. Sie schmeckte nach Sellerie und nach etwas,
das er nicht benennen konnte. Nach Zuhause.

„Du warst lange weg", sagte seine Mutter.

„Ja", sagte Luca.

„Hast du gefunden, was du gesucht hast?"

Luca sah auf den Teller. Die Suppe dampfte. Der Dampf stieg in einer
geraden Linie auf, bevor er sich im Zug der Raumluft auflöste.

„Nein", sagte er. „Und ja."

Seine Mutter sagte nichts. Sie wartete.

„Ich habe zwei Dinge gefunden", sagte Luca. „Die sich nicht vereinbaren
lassen. Und ich habe verstanden, dass ich sie nicht vereinbaren muss."

Er sah auf. Seine Mutter sah ihn an. Ihr Gesicht war alt, aber ihre
Augen waren hell.

„Das ist mehr, als die meisten finden", sagte sie.

„Vielleicht", sagte Luca. „Vielleicht ist es auch weniger."

Er aß weiter. Die Suppe wurde kalt. Die Nacht zog über Boston. Die
Lichter der Stadt brannten. Die Menschen schliefen. Die Algorithmen
rechneten. Das ARS rechnete nicht mehr.

\subsection{IV. Der Rand}\label{iv.-der-rand}

Später, in seinem Zimmer, saß Luca am Schreibtisch. Das Oktavheft lag
vor ihm. Die Seiten waren leer. Er hatte es aus der Quertasche genommen
und aufgeschlagen, ohne zu wissen, warum.

Er nahm einen Stift. Er setzte ihn auf die erste Seite. Er schrieb
nichts. Er wartete.

Und dann schrieb er einen Satz. Einen einzigen Satz. Mit kleiner,
steiler Schrift, so wie er in der Wüste geschrieben hatte, als das Heft
noch voll gewesen war.

Es gibt nicht eine Welt. Es gibt nicht viele Welten. Es gibt den Rand,
an dem beide gelten.

Er legte den Stift hin. Er las den Satz zweimal. Dann schloss er das
Heft.

Er wusste, dass der Satz nicht bleiben würde. Morgen würde er
verschwunden sein. Wie alles andere. Wie Phillips. Wie Martina. Wie die
Parameter, die das ARS überschrieb, wenn man die Senke verließ.

Aber in diesem Moment war er da. Er war geschrieben. Er war gelesen. Er
war gesehen worden.

Und das reichte.

Luca legte sich auf das Bett. Die Decke war dünn. Die Straße war still.
Die Sterne über Boston waren nicht zu sehen, weil die Lichter der Stadt
zu hell waren. Aber Luca wusste, dass sie da waren. Sie waren immer da
gewesen. Sie sahen das Verschwinden. Aber sie bezeugten es nicht.

Er schloss die Augen. Er dachte an Phillips. An den alten Mann in der
Wüste, dessen Hände auf den Oberschenkeln lagen, als wären sie aus
Basalt. An den jüngeren Mann im Trümmerfeld, dessen Jacke in das Grau
der Steine überging. An die beiden Perspektiven, die gleichzeitig
galten.

Er dachte an Martina. An ihr Oktavheft. An den Satz, der nicht
verschwunden war.

Er dachte an den Rand. An den Grenzwert. An das Aufhören, sich zu
entscheiden.

Und dann schlief er ein. Nicht mit einer Antwort. Nicht mit einer Frage.
Nur mit dem Wissen, dass beides da war. Das Ende und die Fülle. Die
Steine und die Möglichkeiten.

Das reichte. Für jetzt.

\section{Über dieses Buch}\label{uxfcber-dieses-buch}

\subsection{Inhalt}\label{inhalt}

Der 37-jährige Doktorand Luca reist in die kalifornische Bimssteinsenke,
um den Jesuiten und Physiker Michael Phillips aufzusuchen. Phillips,
einst Ordinarius an der Gregoriana und Systemarchitekt der
Dialoggrammatiken für InSIM, war maßgeblich an der Entwicklung des
KI/Quanten-Agenten ARS beteiligt, jenem Modell, das Bewusstseinszustände
in iterativen Rekonstruktionen des pompejanischen Sektors V simuliert.
Seit dem plötzlichen Abbruch des Pompeji-Projekts lebt der 78-Jährige
isoliert an einem verlassenen Versuchsfeld in der Wüste. Er hat die
Sprache weitgehend abgelegt; er sitzt auf der Veranda einer Holzbaracke
und fixiert die thermische Erstarrung des Gesteins.

Luca ist der Sohn von Michael Phillips und Maria, einer ehemaligen
Sekretärin am Pontificium Collegium Germanicum et Hungaricum in Rom.
Maria hat Luca allein großgezogen und ihm seine Herkunft nie offenbart.
Sie ist mit ihm nach Boston gegangen, wo Luca promoviert, nicht in
Theologie oder Physik, sondern in einem Grenzbereich zwischen
Sprachmodellierung und Quanteninformatik. Luca weiß, wer sein Vater ist;
die Ähnlichkeit, die Briefe in Marias Schublade und die Spur der
Dialoggrammatiken haben ihm die Gewissheit gegeben, lange bevor er die
Bimssteinsenke erreicht.

Ein anonymer Brief erreicht Luca in Neapel. Kein Absender, keine
Absichtsbezeugung, nur die Koordinaten einer Pufferzone und ein Satz: Er
verweilt im ungemessenen Zustand. Luca reist ab. Er nimmt das vergilbte
Repertoire-Foto des Konsortiums von 1994, mietet einen Wagen in Los
Angeles und bricht in das Versuchsfeld auf.

Am vierzehnten Tag wird das Schweigen der beiden Männer durch die
Ankunft von Martina Rossi unterbrochen. Die 45-jährige Archäologin ist
die Tochter von Michael Phillips und Julia Rossi. Julia hat Martina in
Pompeji allein großgezogen, nachdem sie sich in Rom von Michael getrennt
hatte. Martina tritt ohne Ankündigung in den Raum. Sie berichtet
flüchtig von ihrer Mutter Julia, von archäometrischen Befunden in
Herculaneum und von der Notwendigkeit, sich der Fixierung eines Kollegen
zu entziehen. Luca beobachtet ihre Anwesenheit im isolierten Raum der
Senke, das Fehlen einer ausgeprägten Schattenkontur im späten
Sonnenlicht, die Stille zwischen den Bewegungen.

Am achtzehnten Tag entzieht sich Martina der Messung. Sie hinterlässt
keine Spuren im Schotter; lediglich ihr Oktavheft liegt aufgeschlagen
auf der Tischplatte. Die Seiten sind weiß, bis auf einen einzigen
Eintrag im hinteren Register: In einer anderen Iteration habe ich die
Parameter nicht bestätigt.

Am neunten Tag jedoch, vor Martinas Verschwinden, mitten in die Stille
der Wüste hinein, tritt eine zweite Gestalt in die Senke. Eine
Rekonstruktion des jüngeren Phillips, in der dunklen Baumwolljacke des
Konsortiums von 1994, die behauptet, der Zustand vor der Entscheidung zu
sein. Zwischen dem alten und dem jüngeren Phillips entfaltet sich ein
Gespräch, in dem sich zwei Perspektiven auf denselben Grenzwert
gegenüberstehen: die DeLillo-Perspektive, das Ende, die Steine, die
Entropie, und die Van-Dormael-Perspektive, die Fülle, die Möglichkeiten,
die Differenz. Beide sind richtig. Beide sind nicht vereinbar. Luca
steht zwischen ihnen und begreift, dass der Omegapunkt nicht das eine
oder das andere ist, sondern der Rand, an dem beide gelten, ohne
Synthese, ohne Versöhnung, ohne Ausweg.

Das Verschwinden Martinas destabilisiert das System. Phillips zieht sich
vollkommen in die Rekonstruktionsschleifen des ARS zurück. Luca sucht
das Umland ab, fährt die ausgebleichten Asphaltbänder des Highway 190
ab, befragt die Gestalten an den Tankstellen der Pufferzone und kehrt
ohne Ergebnis zurück. Er beginnt, den Zerfall im Oktavheft zu
protokollieren, die Verdichtung der Zeit, die Quanten-Schnittstellen,
seine Umdeutung von Teilhard de Chardins Omega-Punkt als Grenzzustand
maximaler Entropie.

In den Stunden der maximalen Abkühlung tritt eine Erscheinung an die
Veranda: eine Rekonstruktion des dreißigjährigen Phillips in der Jacke
des Konsortiums von Neapel. Ob es sich um eine Parallelversion Michaels
aus einer anderen Zeitlinie handelt, um eine vom ARS erzeugte
Ersatzfigur oder um Luca selbst, der in einer anderen Iteration als
jüngerer Michael erscheint, bleibt bewusst offen. Das ARS bietet ein
Ersatzsignal an, um den Beobachterstatus aufrechtzuerhalten. Die Entität
artikuliert die Logik der Verschränkung: Martina ist nicht gestorben;
sie hat lediglich den Kohärenzzustand mit dieser spezifischen
Speichereinheit aufgegeben, um in einer parallelen Iteration des Systems
zu verbleiben.

Luca stellt die Messung ein und verlässt die Senke. Als er in Los
Angeles das Flugzeug besteigt und das Notizbuch öffnet, sind alle
Eintragungen gelöscht, das Papier ist makellos und unbeschrieben. In der
Wüste bleibt Phillips zurück. Seine Existenz geht stufenlos in die
reglose Dichte des Bimssteins über.

Nach dem Verschwinden Martinas und dem Erlöschen von Phillips kehrt Luca
nach Boston zurück. Er besucht erneut den Seminarraum, in dem alles
begonnen hat, sieht erneut Mr. Nobody, und führt mit dem Dozenten ein
Gespräch, in dem er die Schlussfolgerung zieht, die die Erzählung
vorbereitet hat: dass DeLillo und Van Dormael sich nicht vereinbaren
lassen, dass sie zwei Perspektiven auf denselben Grenzwert sind, und
dass der einzige ehrliche Umgang mit dem Omegapunkt darin besteht, beide
Perspektiven auszuhalten. Der Epilog führt diese Bewegung in vier Szenen
zu Ende: die Wüste (Phillips\textquotesingle{} Ende), Boston (das Kino),
die Straße (die Brücke) und der Rand (der Satz, der geschrieben wird,
und morgen verschwunden sein wird).

Anmerkung zur Zeitlinie: Die Zeitlosigkeit der Steine spielt nicht in
derselben Zeitlinie wie die übrigen Bände des Pompeji-Projekts. Sie ist
eine Iteration unter vielen, eine mögliche Geschichte im Multiversum, in
der Michael Phillips nicht in Budapest, sondern in der kalifornischen
Wüste endet, in der Martina nicht flieht, sondern sich der Messung
entzieht, und in der Luca nicht als Retter, sondern als Beobachter
auftritt. Die Altersangaben (Luca 37, Martina 45, Michael 78) sind daher
nicht mit den Altersangaben der Trilogie zur Deckung zu bringen, sondern
markieren die spezifische Zeitachse dieser Iteration. Wer die Trilogie
als Multiversumsperspektive liest, wird in dieser Erzählung die
Innenperspektive einer einzelnen Geschichte erkennen, die sich nicht um
chronologische Konsistenz mit den anderen Bänden bemüht, sondern um die
Verdichtung eines einzigen Gedankens: dass am Ende nur das ungemessene
Material bleibt.

\subsection{Hintergrund \& Poetologie}\label{hintergrund-poetologie}

Den Ausgangspunkt für Die Zeitlosigkeit der Steine bildet Don DeLillos
Spätroman Point Omega (2010), der seinerseits auf DeLillos
Auseinandersetzung mit Douglas Gordons Videoinstallation 24 Hour Psycho
im Museum of Modern Art zurückgeht. Die Erzählung greift dabei DeLillos
Interesse an der extrem verlangsamten Wahrnehmung auf und überträgt das
Motiv der zeitlichen Dehnung in den Kontext der modernen
Quantenarchäologie und KI-Theorie.

Das Dreiecksverhältnis aus Point Omega, bei DeLillo vertreten durch den
Intellektuellen Richard Elster, den Filmemacher Jim Finley und Elsters
Tochter Jessie, erfährt in der Erzählung eine grundlegende motivische
Transformation:

- Vom Museum zum Versuchsfeld: An die Stelle der Videokunst tritt die
virtuelle Rekonstruktion von Pompeji durch das Konsortium InSIM. Die
Verlangsamung der Filmkader wird zur Absenkung der Abtastfrequenz eines
Quantenprozessors (ARS).

- Vom Intellektuellen zum Jesuiten-Physiker: Die Rolle der
philosophischen Reflexionsfigur wandelt sich von Richard Elster zu
Michael Phillips. Seine Überlegungen entspringen der Synthese aus
jesuitischer Eschatologie (Pierre Teilhard de Chardin) und der
Everett-Deutsch-Interpretation der Quantenmechanik
(Viele-Welten-Theorie).

- Vom Verschwinden zur Dekohärenz: Das spurlose Verschwinden der Tochter
ist hier kein kriminalistischer oder psychologischer Twist, sondern der
Austritt eines Bewusstseins aus der beobachtbaren Kohärenz einer
kollabierenden Systemiteration.

- Von der Einzelperspektive zur Doppelperspektive: In der Neufassung
wird das Verhältnis zu DeLillo und Van Dormael explizit. Die Erzählung
führt nicht eine Perspektive auf den Omegapunkt vor, sondern zwei, und
macht deren Nicht-Vereinbarkeit sichtbar. Der jüngere Phillips ist keine
DeLillo-Figur. Er ist eine Van-Dormael-Figur: Er spricht von
Möglichkeiten, die nicht verschwinden, von Verzweigungen, die nicht
ineinander aufgehen, von einem Omegapunkt, der nicht das Ende ist,
sondern die Fülle.

Die Erzählung nutzt die kühle, parataktische Ästhetik des
DeLillo-Spätwerks, um die Wüste nicht als physischen Raum, sondern als
isolierten Beobachtungsraum darzustellen. Sprache fungiert darin nicht
als Beschreibungsmittel, sondern als Seismograph einer erstarrenden
Realität, in der am Ende nur das ungemessene Material verbleibt. Dabei
ist Die Zeitlosigkeit der Steine keine chronologische Fortsetzung,
sondern eine eigene Iteration im Multiversum des Pompeji-Projekts, eine
mögliche Geschichte unter vielen, in der dieselben Figuren unter anderen
Anfangsbedingungen andere Wege gehen und in der die Frage nach Schuld,
Beobachtung und Omega-Punkt eine andere Antwort findet.

\subsection{DeLillo und die Stille}\label{delillo-und-die-stille}

In Der Omega-Punkt führt die Stille nicht in Leere, sondern in eine
eigentümliche Zeitlosigkeit, in der Erkenntnis total und nicht mehr
sagbar ist, vielleicht aber erfahrbar. Die Stille hingegen zeigt das
Verstummen der Deutungen: In einem New Yorker Wohnzimmer fallen die
elektronischen Systeme aus, und mit ihnen verschwinden Sprache,
Gewissheit, Ordnung. Zurück bleibt eine Gegenwart, die nichts erklärt
und nichts erlöst.

Zwei Formen der Stille stehen einander gegenüber: die eine als
metaphysische Verdichtung, die andere als reiner Entzug. Die
Zeitlosigkeit der Steine nimmt beide auf. Die Wüste ist der Ort, an dem
die Stille zur physischen Erfahrung wird, nicht als Idylle, sondern als
Zustand nach dem Lärm der Datenströme. Die Steine sind das, was bleibt,
wenn alle Berechnungen eingestellt werden: kein Denkmal, sondern der
Rest.

Die neue Fassung der Erzählung macht diese Spannung explizit. Wo der
alte Phillips das Ende sieht, die Steine, die nicht zeugen, den
Grundzustand, sieht der jüngere Phillips die Fülle. Beide sind
DeLillo-Figuren, aber sie gehören zu verschiedenen DeLillo-Registern:
der alte Phillips zu Die Stille (Entzug, Leere, Wärmetod), der jüngere
Phillips zu einer Lesart, die über Der Omega-Punkt hinausgeht, weil sie
den metaphysischen Trost nicht mehr braucht. Was in der ursprünglichen
Fassung als eine DeLillo-Perspektive erschien, erweist sich in der
Neufassung als gespalten. Der alte Phillips sieht nur noch das Ende. Der
jüngere Phillips sieht beides. Und Luca steht zwischen ihnen und muss
aushalten, dass beide recht haben, aber nicht gleichzeitig.

\subsection{Van Dormael und die
Möglichkeiten}\label{van-dormael-und-die-muxf6glichkeiten}

Jaco Van Dormaels Mr. Nobody (2009) entfaltet ein Universum, in dem jede
Entscheidung eine neue Realität erzeugt. Zeit ist kein Strom, sondern
ein Feld. Der Protagonist durchlebt alle möglichen Leben gleichzeitig,
und jeder Weg ist der richtige Weg. Die „Engel des Vergessens", die
einem bei der Geburt das Wissen über alle Möglichkeiten nehmen, lassen
sich als Metapher für die Viele-Welten-Interpretation lesen, und für den
Omega-Punkt, der jede Version gelebt haben muss, um dieses Wissen zu
haben.

In Das brandneue Testament (2015) wird ein Gottesbild entmachtet, das
aus Angst, Schuld und Kontrolle besteht. Éa, die Tochter Gottes, sendet
allen Menschen ihr Sterbedatum und entzieht dem Vater damit sein
wesentliches Machtinstrument. Nicht durch Widerlegung, sondern durch
Bedeutungslosigkeit verschwindet Gott. Sinn entsteht hier nicht aus
Ordnung, sondern aus Beziehung.

Der fiktive Film Gott wohnt im dritten Stock (aus Drei Filme Zwei
Bücher, Rezension eines fiktiven van Dormael-Filmes neben Das brandneue
Testament und Mr. Nobody als Teil des Pompeji-Projektes) treibt diese
Bewegung weiter. Augustinus, der Name ist Programm, regiert seine
Familie durch ein internalisiertes Regime aus Schuld und Gewissheit.
Seine Tochter Éa verlässt das System und findet unter Obdachlosen eine
fragile Form von Gemeinschaft, die nichts rechtfertigen muss. Am Ende
steht keine Erlösung, sondern eine Verschiebung der Macht: Gott
verschwindet nicht durch Widerlegung, sondern durch Irrelevanz.

Die neue Fassung der Erzählung macht Van Dormael zur zweiten Stimme im
selben Gespräch. Der jüngere Phillips, die Rekonstruktion, die im
Trümmerfeld erscheint, ist keine DeLillo-Figur. Er ist eine
Van-Dormael-Figur: Er spricht von Möglichkeiten, die nicht verschwinden,
von Verzweigungen, die nicht ineinander aufgehen, von einem Omegapunkt,
der nicht das Ende ist, sondern die Fülle. Und er widerspricht dem alten
Phillips, ohne ihn zu widerlegen. Die Erzählung endet nicht bei DeLillo.
Sie endet bei der Frage, ob beide Perspektiven gleichzeitig gelten
können, und bei dem Versuch, sie auszuhalten.

\subsection{Die verbindende Bewegung}\label{die-verbindende-bewegung}

Was all diese Texte verbindet, ist keine Botschaft, sondern eine
Bewegung. Sie führen nicht zu einer neuen Synthese, keinem neuen Bund,
keiner letzten Wahrheit. Sie zeigen vielmehr, wie Macht verschwindet:
die Macht der Technik, des allwissenden Erzählers, des strafenden
Gottes, der großen Sinnnarrative. Nicht durch Zerstörung, sondern durch
Irrelevanz.

Zurück bleibt kein Nihilismus. Sondern eine fragile Form von Humanität.
Still, vorläufig, ohne Garantie. Eine Existenz, die Leid, Zeit und
Endlichkeit nicht aufhebt, aber auch nicht mehr rechtfertigen muss.
Vielleicht ist das der Omega-Punkt: nicht die Verdichtung des
Bewusstseins, sondern seine Ausdünnung bis zu dem Punkt, an dem nur noch
das ungemessene Material bleibt, die Steine.

Die neue Fassung der Erzählung nimmt diesen Satz zurück. Sie setzt nicht
mehr auf eine Antwort, die Ausdünnung, die Steine, das ungemessene
Material, sondern auf die Doppeltheit: dass der Omegapunkt nicht das
Ende oder die Fülle ist, sondern der Rand, an dem beide gelten. Das ist
keine Abschwächung. Es ist die Konsequenz aus der Einsicht, dass DeLillo
und Van Dormael sich nicht vereinbaren lassen, und dass der Versuch, sie
zu vereinbaren, genau das zerstört, was beide zu sagen haben. Teilhard
de Chardin hat eine dieser Perspektiven, die teleologische, für die
ganze gehalten. Er war Augustiner. Er war Thomist. Er konnte nicht
anders, als die Zeit auf ein Ende hin zu denken. Das war seine Größe.
Und sein Irrtum.

In diesem Sinne ist Die Zeitlosigkeit der Steine keine Nacherzählung von
DeLillos oder Van Dormaels Werken. Sie ist eine eigenständige
literarische Ausarbeitung derselben Denkbewegung, die auch der Trilogie
Das Pompeji-Projekt zugrunde liegt. Sie kann parallel zur Trilogie
gelesen werden, als Innenperspektive einer einzelnen Geschichte im
Multiversum, während die Trilogie die Multiversumsperspektive einnimmt.
Für Leser, die philosophisch tiefer einsteigen wollen, ist sie die
konzentrierteste Fassung der Omega-Punkt-Frage im gesamten Werk.

\end{document}
